\section{The chart-lift theorem}\label{sec:pf:hlf-chart}

\providecommand{\hlfcC}{\check{C}}
\providecommand{\hlfcH}{\check{H}}
\providecommand{\hlfcfib}{\Phi}
\providecommand{\hlfcloc}{\lambda}
\providecommand{\hlfcdl}{\operatorname{dlog}}
\providecommand{\hlfcres}{\operatorname{res}}
\providecommand{\hlfcmult}{\operatorname{m}}

This section proves the \emph{chart-lift theorem} (\cref{pf:hlf-chart:main}), the algebraic statement on which the integral case of
the holomorphic Lefschetz formula for finite \'etale correspondences (\cref{thm:galois:hlf}) rests: in the proof of
\lref{UpperTheory.Cb.holomorphic_lefschetz_finiteEtale_integral} (see \cref{sec:pf:geo}, item (e) of the description of the integral
case), the Hodge cohomology $H^p(Z,\Omega^a)$ of a smooth projective variety $Z$ of dimension $d$ is replaced by the \v{C}ech cohomology of the
affine cover defined by a finite morphism $Z\to\mathbb P^d$, the correspondence $(c_1)_*(c_2)^*$ by operators $u_p$ which lift to the
correspondence at the level of \v{C}ech cocycles, and the chart-lift theorem computes $\sum_p(-1)^p\tr(u_p)$ as a sum of local terms over the
fixed points. Everything is formulated in terms of the function fields $K=\C(Z)$, $K'=\C(Y)$ and of subrings of them.

\subsection*{Idea of the proof}
The proof is the classical argument for Lefschetz formulas through the class of the diagonal, carried out explicitly with \v{C}ech and
Koszul complexes. Let $X$ be the smooth projective variety with function field $K$ which is finite over $\mathbb P^d$ via $x$, and $Y$ the
analogous variety for $K'$, with the two maps $c_1,c_2\colon Y\to X$ given by the two embeddings of $K$ into $K'$.
\begin{enumerate}[label=(\arabic*),nosep]
\item A \emph{trace functional} $\theta$ on $\Omega^d_K$ (the $\tau$-coefficient with respect to $\hlfcdl t_1\wedge\dots\wedge\hlfcdl t_d$, traced
down to $F=\C(t_1,\dots,t_d)$, followed by the constant term $\rho$ of Laurent polynomials) induces a pairing between \v{C}ech cocycles of
complementary degrees (\cref{pf:hlf-chart:prb}).
\item The diagonal of $X\times X$ is cut out, in each product chart, by the $2\times2$ minors of the homogeneous coordinates and by
finitely many elements $b\otimes1-1\otimes b$ (\cref{pf:hlf-chart:zfam}). A cocycle $c$ of the associated \v{C}ech--Koszul total complex,
normalized at the local (Koszul) level, exists (\cref{pf:hlf-chart:zex}); its \v{C}ech part is a cocycle of the product cover which acts as the
identity on \v{C}ech cohomology (\cref{pf:hlf-chart:zact}). With the K\"unneth decomposition this gives classes $e_k,f_k$ which are dual
with respect to the pairing (\cref{pf:hlf-chart:zid}); hence $\sum_p(-1)^p\tr(u_p)$ is a sum of pairings $\langle e_k,u(f_k)\rangle$.
\item Pulling $c$ back along the graph $(c_1,c_2)$ of the correspondence gives a cocycle on $Y$ whose class, by the lifting property of the
$u_p$, computes this sum (\cref{pf:hlf-chart:lcup,pf:hlf-chart:edgeb}). On $Y$ the Koszul elements cut out the fixed points; comparing,
at each fixed point, the regular sequence $c_2^*t-c_1^*t$ with the regular sequence $c_1^*t$ by the transition matrix (whose value is
$1-A_e$) shows that the pulled-back class equals the class of an explicit cocycle $h\cdot\hlfcdl t_1\wedge\dots\wedge\hlfcdl t_d$, with $h$
taking the value $\hlfcloc_a(A_e)$ at the fixed point $e$ (\cref{pf:hlf-chart:ymain}).
\item The trace functional of $h\cdot\hlfcdl t_1\wedge\dots\wedge\hlfcdl t_d$ is the sum of the values of $h$ over the fibre of $Y$ above the
origin of $\mathbb P^d$ (\cref{pf:hlf-chart:rt1}).
\end{enumerate}
The development formalizes each step from scratch; the many lemmas of homological algebra (Koszul complexes, filtrations of double complexes,
K\"unneth decomposition) and of commutative algebra (local rings at points, trace duality for finite covers) which it proves along the way are
standard, and we state them and give references or short proofs. Some technical lemmas on explicit cocycles (the local diagonal classes in
\cref{pf:hlf-chart:zex,pf:hlf-chart:zact}, the local cocycles on $Y$ in \cref{pf:hlf-chart:valg,pf:hlf-chart:ybig}) are cited from the development
\cite{LeanSource2026} with their statements and the idea of their proofs.

\subsection*{Organization}
\Cref{pf:hlf-chart:ss:setting} fixes notation and states the theorem; \cref{pf:hlf-chart:ss:red} reduces it to a core statement
(\cref{pf:hlf-chart:core}). \Cref{pf:hlf-chart:ss:forms} treats forms, traces and the trace functional; \cref{pf:hlf-chart:ss:cx}
the weighted \v{C}ech complexes and the Koszul lemmas; \cref{pf:hlf-chart:ss:diag} the diagonal; \cref{pf:hlf-chart:ss:graph} the
pull-back to the correspondence and the localization at the fixed points; \cref{pf:hlf-chart:ss:core} assembles the proof of the core
statement. Throughout, $d\ge0$ is the dimension; it is called $\delta$ in \cref{sec:pf:geo}.

\subsection{Setting and statement}\label{pf:hlf-chart:ss:setting}

\subsubsection*{Fields, coordinates and charts}
Let $F,K,K'$ be fields containing $\C$, with $\C$-algebra maps $F\to K\to K'$; we regard $F$ as a subfield of $K$ and of $K'$ and write
$j_1\colon K\to K'$ for the given (structure) map, which is $F$-linear. Let $j_2\colon K\to K'$ be a second $\C$-algebra map. Let
$x=(x_0,\dots,x_d)\in F^{d+1}$ and $[d]=\{0,1,\dots,d\}$. For a subset $I\subseteq[d]$ let
\[
  A_I=\C\bigl[x_m/x_i : i\in I,\ m\in[d]\bigr]\subseteq F,
\]
and let $B_I\subseteq K$ and $B'_I\subseteq K'$ be the integral closures of $A_I$ in $K$ and in $K'$. Thus $A_\emptyset=\C$, and
$A_{\{i\}}$ is the coordinate ring of the $i$-th standard chart of $\mathbb P^d$ when the $x_m/x_0$ are independent. Put $t_m=x_m/x_0$
($1\le m\le d$), so $t_m\in A_{\{0\}}$. For $J\subseteq[d]$ let $B''_J\subseteq K'$ be the integral closure in $K'$ of $j_2(B_J)$, and for
$I,J\subseteq[d]$ let $B'_{I,J}=B'_I\cdot B''_J$ be the $\C$-subalgebra of $K'$ generated by $B'_I$ and $B''_J$.

\begin{definition}[Finite cover]\label{pf:hlf-chart:def:fincover}
Let $L$ be a field containing $F$ (compatibly with $\C$). We say that $x$ \emph{defines a finite cover} of $L$ if
(i) $x_i\ne0$ for all $i$, the elements $t_1,\dots,t_d$ are algebraically independent over $\C$ and generate $F$ as a field over $\C$;
(ii) $L$ is a finite separable extension of $F$; (iii) for every $i\in[d]$ the integral closure of $A_{\{i\}}$ in $L$ is smooth over $\C$.
\formal{GlobalTheory.PD.IsFinCover}
\end{definition}

Geometrically, $F$ is the function field of $\mathbb P^d$ with homogeneous coordinates $x$, $B_{\{i\}}$ is the coordinate ring of the inverse
image $U_i$ of the $i$-th standard chart in the normalization $X$ of $\mathbb P^d$ in $K$, and $B_I$ is that of $U_I=\bigcap_{i\in I}U_i$. We
shall not use this interpretation; all arguments are algebraic.

\subsubsection*{Forms and \v{C}ech complexes}
For a field $L\supseteq\C$ let $\Omega^q_L=\bigwedge^q_L\Omega_{L/\C}$ and $\Omega^\bullet_L=\bigoplus_q\Omega^q_L$ (the exterior algebra
over $L$), with the universal derivation $d\colon L\to\Omega^1_L$. For a $\C$-subalgebra $B\subseteq L$ let
\[
  \Omega^q(B)=\operatorname{span}_\C\{b\,db_1\wedge\dots\wedge db_q : b,b_1,\dots,b_q\in B\}\subseteq\Omega^q_L,
\]
the \emph{regular $q$-forms over $B$}. Ring maps induce maps of exterior algebras: $j_1$ induces the $K$-algebra map
$j_1^*\colon\Omega^\bullet_K\to\Omega^\bullet_{K'}$ and $j_2$ the $\C$-algebra map $j_2^*\colon\Omega^\bullet_K\to\Omega^\bullet_{K'}$, with
$j_i^*(b\,db_1\wedge\dots)=j_i(b)\,dj_i(b_1)\wedge\dots$. When $K'/K$ is finite separable there is a unique $\C$-linear map
$\Tr\colon\Omega^\bullet_{K'}\to\Omega^\bullet_K$ with $\Tr(b\cdot j_1^*\omega)=\Tr_{K'/K}(b)\,\omega$ for $b\in K'$, $\omega\in\Omega^\bullet_K$
(\cref{pf:hlf-chart:trdr}), the \emph{trace of forms}.

Let $\Sigma$ be a finite totally ordered set and $(M_s)_{s\subseteq\Sigma}$ a family of $\C$-subspaces of a vector space $E$. The
\emph{\v{C}ech complex} $\hlfcC^\bullet(M)$ has as $p$-cochains the families $c=(c_s)$ indexed by the subsets $s\subseteq\Sigma$, with
$c_s\in M_s$ if $s\ne\emptyset$ and $|s|=p+1$, and $c_s=0$ otherwise; the differential is
\begin{equation}\label{eq:pf:hlf-chart:cech}
  (\delta c)_s=\sum_{i\in s}(-1)^{\#\{j\in s:\,j<i\}}c_{s\smallsetminus\{i\}}.
\end{equation}
We write $Z^p,B^p,\hlfcH^p$ for cocycles, coboundaries and cohomology. (If $M_s\subseteq M_{s'}$ for $s\subseteq s'$, which will always be the
case, $\delta$ maps cochains to cochains.) We use two instances:
\begin{itemize}[nosep]
\item $\hlfcC^\bullet(B_\bullet,a)$: $\Sigma=[d]$, $M_I=\Omega^a(B_I)\subseteq\Omega^\bullet_K$; similarly $\hlfcC^\bullet(B'_\bullet,a)$;
\item the \emph{double cover complex} $\hlfcC^\bullet(B'_{\bullet\bullet},a)$: $\Sigma=[d]\sqcup[d]$, the disjoint union of two copies of $[d]$
ordered so that the first copy precedes the second; a subset of $\Sigma$ is written $I\sqcup\bar J$ with $I,J\subseteq[d]$, and
$M_{I\sqcup\bar J}=\Omega^a(B'_{I,J})\subseteq\Omega^\bullet_{K'}$.
\end{itemize}
Note that $B'_{I,\emptyset}=B'_I$ and $B'_{\emptyset,J}=B''_J$ (\cref{pf:hlf-chart:rybasic}); the double cover complex is the \v{C}ech
complex of the cover of $Y$ by the $2(d+1)$ open sets $c_1^{-1}U_i$ and $c_2^{-1}U_j$.

\subsubsection*{Points, values and fixed points}
For a $\C$-subalgebra $B$ of a field $L$, a \emph{character} of $B$ is a $\C$-algebra map $e\colon B\to\C$. The \emph{local ring} of $B$ at $e$
is $B_e=\{v\in L:\exists b\in B,\ e(b)\ne0,\ bv\in B\}$; for $v\in L$ and $\gamma\in\C$ we say that \emph{$v$ has the value $\gamma$ at $e$},
written $v(e)=\gamma$, if there is $b\in B$ with $e(b)\ne0$, $bv\in B$ and $e(bv)=e(b)\gamma$. ($B_e$ is a subring, $e$ extends uniquely to a
character of $B_e$, and $v(e)=\gamma$ if and only if $v\in B_e$ and this extension maps $v$ to $\gamma$; \cref{pf:hlf-chart:valbasic}.)

A character $e$ of $B'_{\{i\}}$ is \emph{fixed} if $j_2(b)(e)=e(j_1(b))$ for every $b\in B_{\{i\}}$ (recall $j_1(B_I)\subseteq B'_I$). The
\emph{fibre over the origin} is the set $\hlfcfib$ of characters $e$ of $B'_{\{0\}}$ with $e(t_1)=\dots=e(t_d)=0$. For
$A\in M_d(\C)$ and $a\ge0$ put
\begin{equation}\label{eq:pf:hlf-chart:locterm}
  \hlfcloc_a(A)=\frac{[X^a]\det(1+XA)}{\det(1-A)}\qquad\bigl([X^a]=\text{coefficient of }X^a\bigr),
\end{equation}
so that $\hlfcloc_a(A)=\tr(\bigwedge^aA)/\det(1-A)$ (\cref{pf:hlf-chart:minorsum}).

\begin{theorem}[Chart-lift theorem]\label{pf:hlf-chart:main}
Let $d,a\ge0$, let $F\to K\to K'$, $x$, $j_2$ be as above, and assume:
\begin{enumerate}[label=\textup{(H\arabic*)},nosep]
\item\label{pf:hlf-chart:H1} $x$ defines a finite cover of $K$ and of $K'$;
\item\label{pf:hlf-chart:H2} for every $i\in[d]$ and every character $e$ of $B'_{\{i\}}$ there is $j\in[d]$ with $B''_{\{j\}}\subseteq(B'_{\{i\}})_e$;
\item\label{pf:hlf-chart:H3} the ring $B'_{\{0\}}/(t_1,\dots,t_d)$ is reduced;
\item\label{pf:hlf-chart:H4} for $i\ne0$ no character of $B'_{\{i\}}$ is fixed;
\item\label{pf:hlf-chart:H5} every fixed character of $B'_{\{0\}}$ vanishes on $t_1,\dots,t_d$;
\item\label{pf:hlf-chart:H6} $\mathcal F\subseteq\hlfcfib$ is a finite set and $e\in\mathcal F$ if and only if $e\in\hlfcfib$ is fixed;
\item\label{pf:hlf-chart:H7} for $e\in\mathcal F$ and $b\in B_{\{0\}}$ with $e(j_1b)=0$ there are $c,g_1,\dots,g_d\in B_{\{0\}}$ with $e(j_1c)\ne0$ and
$cb=\sum_lg_lt_l$;
\item\label{pf:hlf-chart:H8} $(A_e)_{e\in\hlfcfib}$ is a family of matrices in $M_d(\C)$ such that for every $e\in\mathcal F$ and $1\le m\le d$
there are $c,g_1,\dots,g_d\in B'_{\{0\}}$ with $e(c)\ne0$, $c\,j_2(t_m)=\sum_lg_lt_l$ and $e(g_l)=e(c)\,(A_e)_{ml}$ for all $l$;
\item\label{pf:hlf-chart:H9} $\det(1-A_e)\ne0$ for $e\in\mathcal F$;
\item\label{pf:hlf-chart:H10} $(u_p)_{p\ge0}$ are $\C$-linear endomorphisms of $\hlfcH^p(B_\bullet,a)$ which \emph{lift to the correspondence}: for
every $p$-cocycle $b$ of $\hlfcC^\bullet(B_\bullet,a)$ there is a $p$-cocycle $\tilde b$ of $\hlfcC^\bullet(B'_{\bullet\bullet},a)$ such that
$\tilde b_{\emptyset\sqcup\bar J}=j_2^*(b_J)$ for all $J$, the family $(\Tr\tilde b_{I\sqcup\emptyset})_I$ is a $p$-cocycle of
$\hlfcC^\bullet(B_\bullet,a)$, and $u_p[b]=[(\Tr\tilde b_{I\sqcup\emptyset})_I]$;
\item\label{pf:hlf-chart:H11} every $\hlfcH^p(B_\bullet,a)$ is finite-dimensional.
\end{enumerate}
Then
\[
  \sum_{p=0}^{d}(-1)^p\tr(u_p)=\sum_{e\in\mathcal F}\hlfcloc_a(A_e).
\]
\formallabel{GlobalTheory.T.exists_chart_lift_of_finite_cover}\formallabel{UpperTheory.Cb.exists_chart_lift_of_finite_cover}
\end{theorem}

\begin{remark}[On the formal statement]\label{pf:hlf-chart:rem:stmt}
(1) Hypotheses \ref{pf:hlf-chart:H7} and \ref{pf:hlf-chart:H11} are not used in the proof. Finite-dimensionality of the
$\hlfcH^p(B_\bullet,a)$, $p\le d$, is in fact a consequence of the other hypotheses (\cref{pf:hlf-chart:alttrace}). In the development the trace of
an endomorphism of an infinite-dimensional space is defined to be $0$; this convention plays no role here.
(2) In \ref{pf:hlf-chart:H8} the matrix $A_e$ is the matrix of $c_2^*$ on the cotangent space at the fixed point $e$ in the basis
$dt_1,\dots,dt_d$; it is required only for $e\in\mathcal F$, and the values $A_e$ for $e\in\hlfcfib\smallsetminus\mathcal F$ play no role.
(3) Separability in \cref{pf:hlf-chart:def:fincover} is automatic in characteristic $0$; it is part of the formal definition.
(4) The operators $u_p$ enter only through \ref{pf:hlf-chart:H10}, which determines $u_p$ on every class.
\end{remark}

\subsection{Reduction to the core statement}\label{pf:hlf-chart:ss:red}

\begin{lemma}[Forms of high degree]\label{pf:hlf-chart:dlog}
Assume that $x$ defines a finite cover of $L$. Then $\hlfcdl t_m=t_m^{-1}dt_m$ $(1\le m\le d)$ is a basis of the $L$-vector space
$\Omega^1_L$, hence $\Omega^q_L=0$ for $q>d$ and $\Omega^d_L=L\cdot\omega_0$, where
\[
  \omega_0=\hlfcdl t_1\wedge\dots\wedge\hlfcdl t_d .
\]
There is an $L$-linear map $\tau\colon\Omega^\bullet_L\to L$ with $\tau(\omega_0)=1$ and $\tau(\Omega^q_L)=0$ for $q\ne d$, and for every such $\tau$
one has $y=\tau(y)\,\omega_0$ for $y\in\Omega^d_L$.
\formallabel{GlobalTheory.T.dlog_basis_and_vanishing_above_dimension}\formallabel{GlobalTheory.T.exists_topDegreeCoefficientForm_and_eq_on_top_forms}
\end{lemma}

\begin{proof}
$F=\C(t_1,\dots,t_d)$ is purely transcendental, so $dt_1,\dots,dt_d$ is a basis of $\Omega^1_F$, and for the finite separable extension
$L/F$ the natural map $L\otimes_F\Omega^1_F\to\Omega^1_L$ is an isomorphism \cite[\S\S25--26]{Matsumura1989}. As $t_m\ne0$,
the $\hlfcdl t_m$ form a basis as well, and the statements on $\Omega^q_L$ follow. Then $\Omega^\bullet_L=\bigoplus_q\Omega^q_L$ with
$\Omega^d_L$ free of rank one on $\omega_0$; let $\tau$ be the projection to $\Omega^d_L$ followed by the coordinate with respect to
$\omega_0$. If $\tau$ is any map with the stated properties and $y=\gamma\omega_0\in\Omega^d_L$, then $\tau(y)=\gamma$.
\end{proof}

We fix such a $\tau$ for $K$ (the \emph{top coefficient}); for $K'$ we use the analogous $\tau'$. A family $x$ defining a finite cover has
$x_i\ne0$, so the $x_i$ are units of $F$.

The \emph{tangent data} at the fixed points are the following reformulation of \ref{pf:hlf-chart:H8}: for every $e\in\mathcal F$ there are
$g_{ml}\in K'$ ($1\le m,l\le d$) with
\begin{equation}\label{eq:pf:hlf-chart:tan}
  j_2(t_m)=\sum_{l}g_{ml}\,t_l\quad\text{and}\quad g_{ml}(e)=(A_e)_{ml}\quad\text{for all }m,l .
\end{equation}
Indeed, given $c,g_l$ as in \ref{pf:hlf-chart:H8}, we have $c\ne0$ because $e(c)\ne0$, and $g_{ml}=g_l/c$ satisfies
$j_2(t_m)=\sum_lg_{ml}t_l$; and $g_l/c$ has the value $e(g_l)/e(c)=(A_e)_{ml}$ at $e$, since $c\cdot(g_l/c)=g_l\in B'_{\{0\}}$ and
$e(c\cdot g_l/c)=e(c)\cdot e(g_l)/e(c)$.

\begin{theorem}[Core statement]\label{pf:hlf-chart:core}
Assume \ref{pf:hlf-chart:H1}--\ref{pf:hlf-chart:H6} and \ref{pf:hlf-chart:H9}, let $(A_e)_{e\in\hlfcfib}$ satisfy \eqref{eq:pf:hlf-chart:tan}
for every $e\in\mathcal F$, let $a\le d$, and let $(u_p)$ satisfy \ref{pf:hlf-chart:H10}. Then
$\sum_{p=0}^{d}(-1)^p\tr(u_p)=\sum_{e\in\mathcal F}\hlfcloc_a(A_e)$.
\formallabel{UpperTheory.Gal.productCup_of_system.fieldLe_of}
\end{theorem}

The proof of \cref{pf:hlf-chart:core} occupies the rest of the section; it is completed in \cref{pf:hlf-chart:ss:core}.

\begin{lemma}[Reduction]\label{pf:hlf-chart:field0}
\Cref{pf:hlf-chart:core} implies \cref{pf:hlf-chart:main}.
\formallabel{UpperTheory.Gal.productCup_of_system.field0_of}
\end{lemma}

\begin{proof}
If $a\le d$, the hypotheses of \cref{pf:hlf-chart:core} hold by the discussion of \eqref{eq:pf:hlf-chart:tan}, and the conclusions agree.
Let $a>d$. Then $\Omega^a_K=0$ by \cref{pf:hlf-chart:dlog}, so all cochain groups of $\hlfcC^\bullet(B_\bullet,a)$ vanish, $\hlfcH^p(B_\bullet,a)=0$
and $\tr(u_p)=0$ for all $p$. On the other side $\det(1+XA)$ is a polynomial of degree $\le d<a$, so $[X^a]\det(1+XA_e)=0$ and every
$\hlfcloc_a(A_e)$ vanishes. Both sides are $0$.
\end{proof}

\begin{proof}[Proof of \cref{pf:hlf-chart:main}]
Combine \cref{pf:hlf-chart:field0} with \cref{pf:hlf-chart:core}, proved in \cref{pf:hlf-chart:ss:core}.
\end{proof}

\subsection{Charts, traces and the trace functional}\label{pf:hlf-chart:ss:forms}

\subsubsection*{Local rings and charts}

\begin{lemma}[Values]\label{pf:hlf-chart:valbasic}
Let $B$ be a $\C$-subalgebra of a field $L$ and $e$ a character of $B$. Then $B\subseteq B_e$, $B_e$ is a subring of $L$, there is a unique
character $e$ of $B_e$ extending $e$, and $v(e)=\gamma$ holds if and only if $v\in B_e$ and $e(v)=\gamma$. If $v(e)=\gamma\ne0$, then $v\ne0$
and $v^{-1}\in B_e$.
\formallabel{UpperTheory.Gal.productCup_of_system.valBasic_of}
\end{lemma}

\begin{proof}
$B_e$ is the localization of $B$ at the prime ideal $\ker e$, viewed inside $L$; the extension is $v\mapsto e(bv)/e(b)$ for any $b$ with
$e(b)\ne0$, $bv\in B$, which does not depend on $b$ (if $b,b'$ are two choices, $e(b')e(bv)=e(b'bv)=e(b)e(b'v)$). If $v(e)=\gamma\ne0$,
with $b$ as in the definition, then $bv\in B$ and $e(bv)\ne0$, so $v^{-1}=b\cdot(bv)^{-1}\in B_e$.
\end{proof}

\begin{lemma}[Charts]\label{pf:hlf-chart:rybasic}
Assume \ref{pf:hlf-chart:H1} and \ref{pf:hlf-chart:H2}. Then for all $I,I',J,J'\subseteq[d]$:
\begin{enumerate}[label=\textup{(\alph*)},nosep]
\item $B_I\subseteq B_{I'}$ and $B'_I\subseteq B'_{I'}$ for $I\subseteq I'$, and $B'_{I,J}\subseteq B'_{I',J'}$ for $I\subseteq I'$, $J\subseteq J'$;
  $B'_{I,\emptyset}=B'_I$;
\item $j_1(B_I)\subseteq B'_I$ and $j_2(B_I)\subseteq B''_I$;
\item if $j\in J$ and $j'\in[d]$, every $y\in B''_{J\cup\{j'\}}$ satisfies $y\,j_2(x_{j'}/x_j)^N\in B''_J$ for some $N\ge0$; similarly, every
  $y\in B_{I\cup\{i'\}}$ satisfies $y\,(x_{i'}/x_i)^N\in B_I$ for $i\in I$ and some $N$.
  \end{enumerate}
\formallabel{UpperTheory.Gal.productCup_of_system.rYBasic_of}
\end{lemma}

\begin{proof}
(a), (b): $A_I\subseteq A_{I'}$, integral closure is monotone, $j_1$ is the identity on $F$ and $j_2(B_I)$ is integral over $j_2(A_I)$.
$B'_{I,\emptyset}=B'_I\cdot B''_\emptyset$ and $B''_\emptyset$, the integral closure of $j_2(\C)=\C$ in $K'$, is contained in $B'_I$.
(c): $A_{J\cup\{j'\}}=A_J[x_j/x_{j'}]$ for $j\in J$, so $A_{J\cup\{j'\}}$ is the localization of $A_J$ at the powers of $s=x_{j'}/x_j\in A_J$;
since integral closure commutes with localization \cite[Prop.~5.12]{AtiyahMacdonald1969}, $B_{J\cup\{j'\}}=B_J[s^{-1}]$, and the same
holds after applying $j_2$ and taking integral closures in $K'$.
\end{proof}

\begin{lemma}[Finite generation]\label{pf:hlf-chart:ryft}
Assume \ref{pf:hlf-chart:H1}. For $I\ne\emptyset$ the $\C$-algebras $B_I$, $B'_I$ and $B'_{I,J}$ are finitely generated.
\formallabel{UpperTheory.Gal.productCup_of_system.rYFT_of}
\end{lemma}

\begin{proof}
$A_{\{i\}}$ is a polynomial ring and $K'/F$ is finite, so $B'_{\{i\}}$ is a finite $A_{\{i\}}$-module \cite[Cor.~13.13]{Eisenbud1995}, and
$B'_I$ is a localization of $B'_{\{i\}}$ ($i\in I$) at one element by \cref{pf:hlf-chart:rybasic}(c); the same holds for $B_I$. For $J\ne\emptyset$
the same argument applies to $B''_J$, the integral closure of $j_2(A_J)$ in $K'$: $K'$ is finite over $j_2(F)$, being finitely generated over
$\C$ of transcendence degree $d$ and algebraic over $j_2(F)$, which also has transcendence degree $d$. Finally $B''_\emptyset\subseteq B'_I$, so
$B'_{I,J}$ is generated by finitely many elements.
\end{proof}

\begin{lemma}[Regular forms]\label{pf:hlf-chart:omloc}
Let $B\subseteq B'$ be $\C$-subalgebras of a field $L$ and $q\ge0$. Then $\Omega^q(B)\subseteq\Omega^q(B')$, $\Omega^q(B)$ is stable under
multiplication by $B$, $\Omega^p(B)\wedge\Omega^q(B)\subseteq\Omega^{p+q}(B)$, and for every character $e$ of $B$ the span of
$\{\omega:\exists b\in B,\ e(b)\ne0,\ b\omega\in\Omega^q(B)\}$ is $\Omega^q(B_e)$.
\formallabel{UpperTheory.Gal.productCup_of_system.omLoc_of}
\end{lemma}

\begin{proof}
The first three statements are clear from the definition. If $b\omega\in\Omega^q(B)$ with $e(b)\ne0$, then $\omega=b^{-1}(b\omega)$ is a sum of
terms $b^{-1}c\,dc_1\wedge\dots$ with $c,c_i\in B$, and $b^{-1}c\in B_e$; conversely, for $v,v_1,\dots,v_q\in B_e$ choose $b\in B$ with
$e(b)\ne0$ and $bv,bv_i\in B$; then $b^{q+2}v\,dv_1\wedge\dots\wedge dv_q$ lies in $\Omega^q(B)$, as one sees by writing
$b\,dv_i=d(bv_i)-v_i\,db$ and expanding.
\end{proof}

\subsubsection*{The trace of forms}

\begin{lemma}[Trace of forms]\label{pf:hlf-chart:trdr}
Let $K'/K$ be a finite separable extension of fields containing $\C$, with structure map $j_1$. There is a unique $\C$-linear map
$\Tr\colon\Omega^\bullet_{K'}\to\Omega^\bullet_K$ with $\Tr(b\cdot j_1^*\omega)=\Tr_{K'/K}(b)\,\omega$ for $b\in K'$, $\omega\in\Omega^\bullet_K$.
It satisfies the projection formula $\Tr(j_1^*w\wedge y)=w\wedge\Tr(y)$.
\formallabel{GlobalTheory.T.deRhamFormTrace_exists_unique}\formallabel{GlobalTheory.T.deRhamFormTrace_projection_formula}
\end{lemma}

\begin{proof}
$K'\otimes_K\Omega^1_K\to\Omega^1_{K'}$ is an isomorphism (as in the proof of \cref{pf:hlf-chart:dlog}), hence so is
$K'\otimes_K\Omega^\bullet_K\to\Omega^\bullet_{K'}$, $b\otimes\omega\mapsto b\,j_1^*\omega$. Put $\Tr=\Tr_{K'/K}\otimes\mathrm{id}$. Uniqueness holds because
the elements $b\,j_1^*\omega$ span $\Omega^\bullet_{K'}$. For the projection formula take $y=b\,j_1^*\omega$:
$\Tr(j_1^*w\wedge b\,j_1^*\omega)=\Tr(b\,j_1^*(w\wedge\omega))=\Tr_{K'/K}(b)\,w\wedge\omega=w\wedge\Tr(y)$.
\end{proof}

\subsubsection*{The residue functional and the trace functional}
For $\emptyset\ne I\subseteq[d]$ let
\[
  \Omega_F(I)=\Bigl\{f\in F : f\cdot\prod_{l\in[d]}\frac{x_i}{x_l}\in A_I\ \text{for all }i\in I\Bigr\}.
\]
On the chart $x_i\ne0$ of $\mathbb P^d$, with coordinates $s_l=x_l/x_i$ ($l\ne i$), one has
$\hlfcdl t_1\wedge\dots\wedge\hlfcdl t_d=\pm\prod_{l\ne i}s_l^{-1}\,\bigwedge_{l\ne i}ds_l$; so $f\in\Omega_F(I)$ means that the $d$-form
$f\omega_0$ of $F$ is regular on $\bigcap_{i\in I}\{x_i\ne0\}$.

\begin{lemma}[Residue functional]\label{pf:hlf-chart:rho}
Assume that $t_1,\dots,t_d$ are algebraically independent over $\C$ and generate $F$. There is a $\C$-linear map $\rho\colon F\to\C$ with
$\rho(1)=1$ and $\rho(\prod_jx_j^{m_j})=0$ for every $0\ne m\in\Z^{d+1}$ with $\sum_jm_j=0$. Every such $\rho$ satisfies:
\begin{enumerate}[label=\textup{(\alph*)},nosep]
\item $\rho(f)=f(0)$ for $f\in A_{\{0\}}=\C[t_1,\dots,t_d]$;
\item $\rho(\Omega_F(I))=0$ for every nonempty $I\subsetneq[d]$.
\end{enumerate}
\formallabel{GlobalTheory.T.exists_laurentConstantTermForm_functionField}\formallabel{GlobalTheory.T.residue_functional_vanishes_on_omega_forms}
\end{lemma}

\begin{proof}
The monomials $\prod_jx_j^{m_j}$ with $\sum m_j=0$ are the Laurent monomials $\prod_{l\ge1}t_l^{m_l}$; they are linearly independent over $\C$
by the algebraic independence of the $t_l$, and $m\mapsto\prod_jx_j^{m_j}$ is injective on $\{\sum m_j=0\}$. Define $\rho$ on their span by
$\rho(1)=1$ and $\rho=0$ on the other monomials, and extend it linearly to $F$ (choose a complement of the span). (a) is immediate since
$t^n=\prod_jx_j^{m_j}$ with $m_0=-\sum n_l$, $m_l=n_l$. (b) Let $i\in I$ and $l_0\notin I$. An element of $A_I$ is a linear combination of
monomials $\prod_jx_j^{m_j}$ with $\sum m_j=0$ and $m_j\ge0$ for $j\notin I$, since the generators $x_m/x_{i'}$, $i'\in I$, have negative
exponents only at indices in $I$. If $f\in\Omega_F(I)$, then $f=g\prod_l(x_l/x_i)$ with $g\in A_I$, and $\prod_l(x_l/x_i)$ has exponent $1$ at
every $l\ne i$. So $f$ is a combination of monomials whose exponent at $l_0$ is $\ge1$; these monomials are $\ne1$, and $\rho(f)=0$.
\end{proof}

We fix such a $\rho$. The \emph{trace functional} of $K$ is the $\C$-linear map
\[
  \theta=\rho\circ\Tr_{K/F}\circ\tau\colon\Omega^\bullet_K\to\C,
\]
and similarly $\theta'=\rho\circ\Tr_{K'/F}\circ\tau'$ for $K'$. By \cref{pf:hlf-chart:dlog}, $\theta$ vanishes off degree $d$ and $\theta(h\omega_0)=\rho(\Tr_{K/F}h)$.

\begin{littheorem}[Trace duality for finite covers; formalized]\label{pf:hlf-chart:dual}
Assume that $x$ defines a finite cover of $L$, let $\tau_L$ be a top coefficient (\cref{pf:hlf-chart:dlog}), $B_I(L)$ the integral closure of
$A_I$ in $L$, and $\emptyset\ne I\subseteq[d]$.
\begin{enumerate}[label=\textup{(\alph*)},nosep]
\item For $z\in L$: $\Tr_{L/F}(zb)\in\Omega_F(I)$ for all $b\in B_I(L)$ if and only if $z=\tau_L(y)$ for some $y\in\Omega^d(B_I(L))$.
\item For $p+q=d$ and $y\in\Omega^q_L$: $y\in\Omega^q(B_I(L))$ if and only if $\Tr_{L/F}\tau_L(w\wedge y)\in\Omega_F(I)$ for all $w\in\Omega^p(B_I(L))$.
\end{enumerate}
\formallabel{GlobalTheory.T.complementaryModule_standardOpens_iff}\formallabel{GlobalTheory.T.twistedDual_regularForms_finiteCover}
\end{littheorem}

In geometric terms: let $\pi\colon U\to V$ be the finite morphism from $U=\operatorname{Spec}B_I(L)$ (smooth) to the open subset
$V=\operatorname{Spec}A_I$ of $\mathbb P^d$. Statement (a) says that the trace of differential forms identifies $\pi_*\Omega^d_U$ with
$\mathcal Hom_{\mathcal O_V}(\pi_*\mathcal O_U,\Omega^d_V)$, i.e.\ $\pi^!\Omega^d_V=\Omega^d_U$: this is duality for finite morphisms of smooth
varieties \cite[III, Ex.~6.10 and Ex.~7.2]{Hartshorne1977}, in the explicit form in which the abstract trace map is the trace of differential
forms and the dual of $\pi_*\mathcal O_U$ is Dedekind's complementary module (see \cite{Kunz1986} and \cite{Lipman1984}). Statement (b)
follows from (a) and the perfectness of the pairing $\Omega^p_U\times\Omega^q_U\to\Omega^d_U$, $\Omega^1_U$ being locally free of rank $d$.
The development proves both statements directly, through the complementary module $\{z:\Tr_{L/F}(zB_I(L))\subseteq\Omega_F(I)\}$; we use
them as stated.

\begin{lemma}[Traces of regular top forms]\label{pf:hlf-chart:trreg}
Assume \ref{pf:hlf-chart:H1}. Let $I\subseteq[d]$ be nonempty.
\begin{enumerate}[label=\textup{(\alph*)},nosep]
\item If $y\in\Omega^d(B_I)$, then $\Tr_{K/F}\tau(y)\in\Omega_F(I)$; if moreover $I\ne[d]$, then $\theta(y)=0$.
\item If $y\in\Omega^d(B'_I)$, then $\Tr(y)\in\Omega^d(B_I)$; if moreover $I\ne[d]$, then $\theta(\Tr y)=0$.
\end{enumerate}
\end{lemma}

\begin{proof}
(a) Apply \cref{pf:hlf-chart:dual}(a) with $z=\tau(y)$ and $b=1$, then \cref{pf:hlf-chart:rho}(b). (b) Let $z'=\tau'(y)$, so $y=z'\omega'_0$
with $\omega'_0=j_1^*\omega_0$ (the $t_m$ lie in $F$), and $\Tr(y)=\Tr_{K'/K}(z')\,\omega_0$ by \cref{pf:hlf-chart:trdr}. For $b\in B_I$,
$\Tr_{K/F}(\Tr_{K'/K}(z')\,b)=\Tr_{K'/F}(z'j_1(b))\in\Omega_F(I)$ by \cref{pf:hlf-chart:dual}(a) for $K'$, since $j_1(b)\in B'_I$. By
\cref{pf:hlf-chart:dual}(a) for $K$, $\Tr_{K'/K}(z')=\tau(y_1)$ with $y_1\in\Omega^d(B_I)$, and $\Tr(y)=\tau(y_1)\omega_0=y_1$. The last
assertion follows from (a).
\end{proof}

\subsubsection*{The pairing}
For families $e=(e_s)$ and $f=(f_s)$ indexed by subsets of $[d]$ (with values in $\Omega^\bullet_K$) put
\begin{equation}\label{eq:pf:hlf-chart:pair}
  \langle e,f\rangle=\sum_{j=0}^{d}\theta\bigl(e_{\{0,\dots,j\}}\wedge f_{\{j,\dots,d\}}\bigr).
\end{equation}
If $e$ is a $(d-q)$-cochain and $f$ a $q$-cochain of \v{C}ech complexes, only $j=d-q$ contributes. More generally, for a bilinear map
$\mu\colon E\times E'\to E''$ the \emph{cup product} of families $e,f$ is $(e\cup_\mu f)_s=\sum_{j\in s}\mu(e_{\{i\in s:i\le j\}},f_{\{i\in s:i\ge j\}})$,
so that $\langle e,f\rangle=\theta((e\cup_\wedge f)_{[d]})$. For \v{C}ech cochains of degrees $p,q$ it satisfies the Leibniz rule
$\delta(e\cup_\mu f)=\delta e\cup_\mu f+(-1)^pe\cup_\mu\delta f$ (a direct computation with \eqref{eq:pf:hlf-chart:cech}).

\begin{lemma}[Adjointness]\label{pf:hlf-chart:adj}
Let $(M_s)$, $(M'_s)$ be increasing families of subspaces (as for \v{C}ech complexes on $[d]$) and $\beta\colon E\times E'\to\C$ bilinear with
$\beta(M_{[d]\smallsetminus\{i\}},M'_{[d]\smallsetminus\{i\}})=0$ whenever $[d]\smallsetminus\{i\}\ne\emptyset$. Then for each $n$ there is a sign $\varepsilon$
such that $\beta$-pairing \eqref{eq:pf:hlf-chart:pair} satisfies $\langle\delta e,f\rangle_\beta=\varepsilon\langle e,\delta f\rangle_\beta$ for all
$n$-cochains $e$ of $\hlfcC^\bullet(M)$ and $(d-n-1)$-cochains $f$ of $\hlfcC^\bullet(M')$.
\formallabel{GlobalTheory.T.cechDifferential_adjoint_cupProductPairing}
\end{lemma}

\begin{proof}
$\langle\delta e,f\rangle_\beta+(-1)^n\langle e,\delta f\rangle_\beta=\beta$ applied to $\delta(e\cup f)_{[d]}=\sum_i\pm(e\cup f)_{[d]\smallsetminus\{i\}}$,
and every $(e\cup f)_{[d]\smallsetminus\{i\}}$ is a sum of products of elements of $M_{[d]\smallsetminus\{i\}}$ and $M'_{[d]\smallsetminus\{i\}}$
(monotonicity), on which $\beta$ vanishes. Take $\varepsilon=-(-1)^n$.
\end{proof}

\begin{lemma}[The pairing descends to cohomology]\label{pf:hlf-chart:prb}
Assume \ref{pf:hlf-chart:H1}, $a\le d$, $q\le d$. If $e\in Z^{d-q}(\hlfcC^\bullet(B_\bullet,d-a))$ and $y\in B^q(\hlfcC^\bullet(B_\bullet,a))$, then
$\langle e,y\rangle=0$.
\formallabel{UpperTheory.Gal.productCup_of_system.prB_of}
\end{lemma}

\begin{proof}
For $i\in[d]$ with $I=[d]\smallsetminus\{i\}\ne\emptyset$, $w\in\Omega^{d-a}(B_I)$ and $w'\in\Omega^a(B_I)$ we have $w\wedge w'\in\Omega^d(B_I)$,
so $\theta(w\wedge w')=0$ by \cref{pf:hlf-chart:trreg}(a). By \cref{pf:hlf-chart:adj}, if $y=\delta y'$ then
$\langle e,\delta y'\rangle=\pm\langle\delta e,y'\rangle=0$.
\end{proof}

\begin{lemma}[Trace and the top degree]\label{pf:hlf-chart:trtop}
Assume \ref{pf:hlf-chart:H1}. Then $\Tr(j_1^*w\wedge y)=w\wedge\Tr(y)$; $\theta(\Tr(h\omega'_0))=\rho(\Tr_{K'/F}h)$ for $h\in K'$; and
$\theta(\Tr(y_{[d]}))=0$ for every $y\in B^d(\hlfcC^\bullet(B'_\bullet,d))$.
\formallabel{UpperTheory.Gal.productCup_of_system.trTop_of}
\end{lemma}

\begin{proof}
The first identity is \cref{pf:hlf-chart:trdr}. Next, $\Tr(h\omega'_0)=\Tr(h\,j_1^*\omega_0)=\Tr_{K'/K}(h)\omega_0$, so
$\theta(\Tr(h\omega'_0))=\rho(\Tr_{K/F}\Tr_{K'/K}h)=\rho(\Tr_{K'/F}h)$. Finally, if $y=\delta w$, then
$y_{[d]}=\sum_i\pm w_{[d]\smallsetminus\{i\}}$ with $w_{[d]\smallsetminus\{i\}}\in\Omega^d(B'_{[d]\smallsetminus\{i\}})$ (and $y=0$ if $d=0$), and each
$\theta(\Tr w_{[d]\smallsetminus\{i\}})$ vanishes by \cref{pf:hlf-chart:trreg}(b).
\end{proof}

\begin{lemma}[Residue of a trace]\label{pf:hlf-chart:rt1}
Let $x$ define a finite cover of $L$ with $B_{\{0\}}(L)/(t_1,\dots,t_d)$ reduced, and let $\hlfcfib_L$ be the set of characters of $B_{\{0\}}(L)$
vanishing on $t_1,\dots,t_d$. Then $\hlfcfib_L$ is finite and for every $h\in B_{\{0\}}(L)$
\[
  \rho\bigl(\Tr_{L/F}h\bigr)=\theta_L(h\omega_0)=\sum_{e\in\hlfcfib_L}e(h).
\]
\formallabel{GlobalTheory.T.trace_cechCochain_fraction_eq_sum_fibre}\formallabel{UpperTheory.Gal.productCup_of_system.deRham_restriction_trace_degree_one}
\end{lemma}

\begin{proof}
We give the standard argument. Let $A=A_{\{0\}}=\C[t_1,\dots,t_d]$, $B=B_{\{0\}}(L)$, $\mathfrak m=(t_1,\dots,t_d)\subseteq A$. $B$ is a finite
$A$-module \cite[Cor.~13.13]{Eisenbud1995}; it is smooth over $\C$, hence regular, hence Cohen--Macaulay, and a finite extension of the regular
ring $A$ by a Cohen--Macaulay ring of the same dimension is flat \cite[Thm.~23.1]{Matsumura1989}. So $B_\mathfrak m=B\otimes_AA_\mathfrak m$ is free of
rank $n=[L:F]$ over the local ring $A_\mathfrak m$, and $\Tr_{L/F}(h)$ is the trace of multiplication by $h$ on $B_\mathfrak m$ (base change to $F$).
Hence $\Tr_{L/F}(h)\in A_\mathfrak m$, and $\Tr_{L/F}(h)\in A$ because $A$ is integrally closed and $\Tr_{L/F}h$ is integral over $A$
\cite[Prop.~5.15]{AtiyahMacdonald1969}. Reducing modulo $\mathfrak m$, $\Tr_{L/F}(h)(0)$ is the trace of multiplication by $h$ on the finite
$\C$-algebra $B/\mathfrak mB$; this algebra is reduced, hence isomorphic to $\C^{\hlfcfib_L}$ via its characters, which are the elements of
$\hlfcfib_L$; so the trace is $\sum_{e\in\hlfcfib_L}e(h)$. Finally $\rho(\Tr_{L/F}h)=\Tr_{L/F}(h)(0)$ by \cref{pf:hlf-chart:rho}(a), and
$\theta_L(h\omega_0)=\rho(\Tr_{L/F}h)$ by definition.
\end{proof}

\subsubsection*{Traces of linear maps}

\begin{lemma}[Traces through dual families]\label{pf:hlf-chart:trlin}
Let $V$ be a $\C$-vector space, $\varphi_1,\dots,\varphi_n\in V^*$ and $v_1,\dots,v_n\in V$ with $y=\sum_k\varphi_k(y)v_k$ for all $y\in V$.
Then $V$ is finite-dimensional and $\tr(u)=\sum_k\varphi_k(u(v_k))$ for every endomorphism $u$ of $V$.
\formallabel{UpperTheory.Gal.productCup_of_system.trLin_of}
\end{lemma}

\begin{proof}
$V$ is spanned by the $v_k$. With $P\colon\C^n\to V$, $P(c)=\sum c_kv_k$, and $Q\colon V\to\C^n$, $Q(y)=(\varphi_k(y))_k$, we have $PQ=\mathrm{id}_V$,
so $\tr(u)=\tr(uPQ)=\tr(QuP)=\sum_k\varphi_k(u(v_k))$.
\end{proof}

\begin{lemma}[Alternating trace through a dual decomposition]\label{pf:hlf-chart:alttrace}
Assume \ref{pf:hlf-chart:H1} and $a\le d$. Let $(p_k,e_k,f_k)_{1\le k\le M}$ be triples with $p_k\le d$,
$e_k\in Z^{d-p_k}(\hlfcC^\bullet(B_\bullet,d-a))$, $f_k\in Z^{p_k}(\hlfcC^\bullet(B_\bullet,a))$, such that for every $q$ and every
$\alpha\in Z^q(\hlfcC^\bullet(B_\bullet,a))$
\begin{equation}\label{eq:pf:hlf-chart:dualdec}
  \alpha\equiv\sum_{k:\,p_k=q}\langle e_k,\alpha\rangle\,f_k\pmod{B^q}.
\end{equation}
Let $u_q$ be endomorphisms of $\hlfcH^q(B_\bullet,a)$ and $g_k\in Z^{p_k}$ with $u_{p_k}[f_k]=[g_k]$. Then every $\hlfcH^q(B_\bullet,a)$ is
finite-dimensional and $\sum_{q=0}^d(-1)^q\tr(u_q)=\sum_k(-1)^{p_k}\langle e_k,g_k\rangle$.
\end{lemma}

\begin{proof}
By \cref{pf:hlf-chart:prb}, $\varphi_k[\alpha]=\langle e_k,\alpha\rangle$ is well defined on $\hlfcH^{p_k}$; by \eqref{eq:pf:hlf-chart:dualdec} the
$\varphi_k,[f_k]$ with $p_k=q$ satisfy the hypothesis of \cref{pf:hlf-chart:trlin} on $V=\hlfcH^q$ (for $q>d$ the sum is empty and $\hlfcH^q=0$).
Hence $\tr(u_q)=\sum_{p_k=q}\langle e_k,g_k\rangle$; multiply by $(-1)^q$ and sum over $q\le d$.
\end{proof}

\begin{lemma}[Principal minors]\label{pf:hlf-chart:minorsum}
For $A\in M_d(\C)$, $[X^a]\det(1+XA)=\sum_{|S|=a}\det(A_{S,S})=\tr\bigl(\bigwedge^aA\bigr)$, the sum running over the subsets $S\subseteq\{1,\dots,d\}$
of cardinality $a$ and $A_{S,S}$ denoting the principal submatrix.
\formallabel{UpperTheory.Gal.productCup_of_system.minorSum_of}
\end{lemma}

\begin{proof}
Expand $\det(1+XA)$ by multilinearity in the columns: choosing the column of $XA$ for $j\in S$ and of $1$ for $j\notin S$ gives
$X^{|S|}\det(A_{S,S})$. The matrix of $\bigwedge^aA$ in the basis $e_S$ has diagonal entries $\det(A_{S,S})$.
\end{proof}

\subsection{Weighted \v{C}ech complexes and Koszul complexes}\label{pf:hlf-chart:ss:cx}

All complexes of the proof are instances of one construction, which treats \v{C}ech differentials and Koszul differentials alike.

\begin{definition}[Weighted \v{C}ech complexes]\label{pf:hlf-chart:def:wcx}
Let $\Sigma$ be a finite totally ordered set, $E$ a $\C$-vector space, $(A_s)_{s\subseteq\Sigma}$ subspaces of $E$ and $w\colon\Sigma\to\End_\C(E)$
pairwise commuting endomorphisms (the \emph{weights}) with $w_i(A_s)\subseteq A_{s\cup\{i\}}$ for $i\notin s$. The complex $\mathcal C(A,w)$ has
$n$-cochains the families $c=(c_s)_{s\subseteq\Sigma}$ with $c_s\in A_s$ if $|s|=n$ and $c_s=0$ otherwise, and differential
\[
  (\partial_wc)_s=\sum_{i\in s}(-1)^{\#\{j\in s:\,j<i\}}\,w_i\bigl(c_{s\smallsetminus\{i\}}\bigr).
\]
\formal{GlobalTheory.WH.acx}
\end{definition}

Thus $\partial_w^2=0$, and $\partial_{w+w'}=\partial_w+\partial_{w'}$, $\partial_w\partial_{w'}+\partial_{w'}\partial_w=0$ if $w,w'$ commute and have
disjoint supports (\lref{UpperTheory.Gal.productCup_of_system.wdBasic_of}, a direct computation). Examples:
\begin{enumerate}[label=(\roman*),nosep]
\item $w\equiv1$ and $A_\emptyset=0$: the \v{C}ech complex $\hlfcC^\bullet(A)$ of \eqref{eq:pf:hlf-chart:cech}, with degrees shifted by one.
\item $A_s=M$ for all $s$ and $w_l$ the action of elements $r_l$ of a commutative ring $S$ acting on $E$, $M$ being stable under a subring $R$
containing the $r_l$: the (cohomological) \emph{Koszul complex} $K(r;M)$. Its top cohomology is $M/(r)M$, where $(r)M=\sum_lr_lM$.
\item $\Sigma=[d]\sqcup[d]$ (first copy first), $w\equiv1$ and $A_{I\sqcup\bar J}=M_{I,J}$ if $I,J\ne\emptyset$, $A_{I\sqcup\bar J}=0$ otherwise:
the \emph{mixed complex} $\mathrm{Mix}(M)$ of a family $(M_{I,J})$; it is the total complex of the \v{C}ech double complex of a product of two
covers, the degree of the $(I,J)$-component being $|I|+|J|$ (\v{C}ech bidegree plus $2$).
\end{enumerate}
For $\Sigma=\Sigma_1\sqcup\Sigma_2$ ($\Sigma_1$ first) and $s=b\sqcup L$, $(\partial_wc)_{b\sqcup L}$ is the sum of the differential in the
$\Sigma_1$-variables at fixed $L$ and $(-1)^{|b|}$ times the differential in the $\Sigma_2$-variables at fixed $b$
(\lref{UpperTheory.Gal.productCup_of_system.wdUni_of}). For a cochain $c$ and $b\subseteq\Sigma_1$, the \emph{column} of $c$ at $b$ is
$\mathrm{col}_b(c)=(c_{b\sqcup L})_{L\subseteq\Sigma_2}$.

\begin{definition}[Systems and \v{C}ech--Koszul complexes]\label{pf:hlf-chart:def:sys}
A \emph{system} $T$ (of dimension $d$) consists of a commutative $\C$-algebra $S$ acting on a $\C$-vector space $E$, units
$X_0,\dots,X_d$ and $Y_0,\dots,Y_d$ of $S$, subalgebras $R_{I,J}\subseteq S$ and subspaces $M_{I,J}\subseteq E$ ($I,J\subseteq[d]$). It is
\emph{good} if, for nonempty $I\subseteq I'$, $J\subseteq J'$: $R_{I,J}\subseteq R_{I',J'}$, $M_{I,J}\subseteq M_{I',J'}$, $R_{I,J}M_{I,J}\subseteq M_{I,J}$,
and $X_m/X_i,\ Y_m/Y_j\in R_{I,J}$ for $i\in I$, $j\in J$, $m\in[d]$. \emph{Koszul data} for $T$ are a finite totally ordered set $\kappa$,
bidegrees $n_l=(n_l',n_l'')\in\N^2$ and elements $z_l\in S$ ($l\in\kappa$); their \emph{dehomogenizations} are
\[
  u_l(i,j)=X_i^{n'_l}Y_j^{n''_l},\qquad r_l(i,j)=z_l/u_l(i,j)\qquad(i,j\in[d]),
\]
and the data are \emph{admissible} if $r_l(i,j)\in R_{\{i\},\{j\}}$ for all $i,j,l$. The \emph{\v{C}ech--Koszul complex}
$\mathrm{Tot}_T(n,z)$ is the weighted \v{C}ech complex on $\Sigma=[d]\sqcup[d]\sqcup\kappa$ (in this order) with
\[
  A_{I\sqcup\bar J\sqcup L}=u_L(\min I,\min J)\,M_{I,J}\ \ (I,J\ne\emptyset),\qquad A_{I\sqcup\bar J\sqcup L}=0\ \text{otherwise},
\]
where $u_L=\prod_{l\in L}u_l$, and weights $1$ on $[d]\sqcup[d]$ and $z_l$ on $l\in\kappa$. Its $(I,J)$-\emph{column complex} is the weighted complex
on $\kappa$ with modules $u_L(\min I,\min J)M_{I,J}$ and weights $z_l$. The \emph{mixed complex} of $T$ is $\mathrm{Mix}(T)=\mathrm{Mix}(M)$, and
$\mathrm{top}(c)=(c_{I\sqcup\bar J\sqcup\emptyset})_{I,J}$.
\formal{UpperTheory.Gal.productCup_of_system.Sys.cx}
\end{definition}

The \emph{diagonal} Koszul elements of a system are the $2\times2$ minors $\mu_{(m,m')}=X_mY_{m'}-X_{m'}Y_m$, indexed by the ordered pairs
$(m,m')\in[d]^2$ (lexicographic order), of bidegree $(1,1)$. For $j\in[d]$ we write $D_j=\{(m,j):m\ne j\}$, a set of $d$ pairs.

\begin{lemma}[Basic properties]\label{pf:hlf-chart:sysbasic}
Let $T$ be a good system with admissible Koszul data $(n,z)$, and $N\in\Z$.
\begin{enumerate}[label=\textup{(\alph*)},nosep]
\item $\mathrm{Tot}_T(n,z)$ is a complex; a family $c$ is an $N$-cochain if and only if each column $\mathrm{col}_{I,J}(c)$ is an
  $(N-|I|-|J|)$-cochain of the $(I,J)$-column complex (and $c$ vanishes where $I$ or $J$ is empty).
\item $\mathrm{top}$ maps cocycles to cocycles and coboundaries to coboundaries of $\mathrm{Mix}(T)$.
\item If $c$ is an $N$-cocycle whose components with $|I'|+|J'|<|I|+|J|$ vanish, then $\mathrm{col}_{I,J}(c)$ is a cocycle of the column complex.
\item The restriction of a cocycle to the index sets $I\sqcup\bar J\sqcup L$ with $L$ contained in a subset $\kappa'\subseteq\kappa$ is a cocycle of
  $\mathrm{Tot}_T(n|_{\kappa'},z|_{\kappa'})$.
\item The $(I,J)$-column complex is isomorphic, by $c\mapsto(u_L^{-1}c_L)_L$, to the Koszul complex $K(r(\min I,\min J);M_{I,J})$, and the
  column complexes are monotone: $Z^m$ of the $(I,J)$-column is contained in $Z^m$ of the $(I',J')$-column for $I\subseteq I'$, $J\subseteq J'$.
\end{enumerate}
\formallabel{UpperTheory.Gal.productCup_of_system.sysBasic_of}\formallabel{UpperTheory.Gal.productCup_of_system.sysOk_of}\formallabel{UpperTheory.Gal.productCup_of_system.rS_of}
\end{lemma}

\begin{proof}
The weights satisfy the condition of \cref{pf:hlf-chart:def:wcx}: with $i=\min I$, $j=\min J$,
\[
  z_l\,u_L(i,j)M_{I,J}=u_{L\cup\{l\}}(i,j)\,r_l(i,j)M_{I,J},
\]
and $r_l(i,j)\in R_{I,J}$ by admissibility and goodness; moreover, when $I$ grows, $\min I$ may change, and the quotient
$u_L(i,j)/u_L(i',j')$ is a product of the elements $X_i/X_{i'}$, $Y_j/Y_{j'}$, which lie in $R_{I',J'}$. This gives (a) and the monotonicity in (e).
(b): the differential of a component with $L=\emptyset$ only involves components with $L=\emptyset$. (c): by the decomposition of $\partial$ along
$\Sigma=([d]\sqcup[d])\sqcup\kappa$, the $(I,J)$-column of $\partial c$ is the column differential of $\mathrm{col}_{I,J}(c)$ plus \v{C}ech terms
involving components with smaller $|I'|+|J'|$. (d): the differential of a component indexed by $s$ only involves components indexed by subsets of
$s$. (e): $z_l=r_lu_l$.
\end{proof}

\subsubsection*{Koszul complexes}
In the following lemmas $S$ acts on $E$, $R\subseteq S$ is a subring, $M\subseteq E$ is an $R$-stable subspace and $r=(r_l)_{l\in\kappa}$ is a
finite family in $R$; we write $(r)M=\sum_lr_lM$. We say that $K(r;M)$ is \emph{exact below $d$} if $Z^m\subseteq B^m$ for all $m<d$. A sequence
$f_1,\dots,f_d\in S$ is \emph{$M$-regular} if $f_km\in(f_1,\dots,f_{k-1})M$ implies $m\in(f_1,\dots,f_{k-1})M$ for $m\in M$ (no
condition $M\ne(f)M$ is imposed).

\begin{lemma}[Koszul complexes]\label{pf:hlf-chart:koszul}
\begin{enumerate}[label=\textup{(\alph*)},nosep]
\item If $g=\sum_la_lr_l$ with $a_l\in R$, then multiplication by $g$ maps $Z^n(K(r;M))$ into $B^n(K(r;M))$.
\item If $f_1,\dots,f_d$ is an $M$-regular sequence with every $f_k\in\sum_lRr_l$, then $K(r;M)$ is exact below $d$.
\item Every coboundary of $K(r;M)$ has all components in $(r)M$. If $e\colon\{1,\dots,d\}\to\kappa$ is an order embedding and $K(r\circ e;M)$ is
  exact below $d$, then so is $K(r;M)$, and:
\item for $m\in M$ with $r_lm\in(r\circ e)M$ for all $l\in\kappa$, there is $c\in Z^d(K(r;M))$ with $c_{e(\{1..d\})}=m$;
\item a cocycle $c\in Z^d(K(r;M))$ with $c_{e(\{1..d\})}\in(r\circ e)M$ is a coboundary;
\item if $e'$ is a second order embedding with image disjoint from that of $e$ and $r_{e'(a)}=\sum_jN_{aj}r_{e(j)}$ with $N_{aj}\in R$, then
  $c_{e'(\{1..d\})}-\det(N)\,c_{e(\{1..d\})}\in(r\circ e)M$ for every $c\in Z^d(K(r;M))$.
\end{enumerate}
\formallabel{UpperTheory.Gal.productCup_of_system.htpy_of}\formallabel{UpperTheory.Gal.productCup_of_system.firstCohomologyBound_of_boundaryBound}\formallabel{UpperTheory.Gal.productCup_of_system.kSub_of}\formallabel{UpperTheory.Gal.productCup_of_system.koszulTwoExact_of_gl2}\formallabel{UpperTheory.Gal.productCup_of_system.koszulTwoInjective_of_gl1}\formallabel{UpperTheory.Gal.productCup_of_system.koszulThree_of_det}
\end{lemma}

\begin{proof}
(a) The map $h(c)_s=\sum_la_l\,(\pm c_{s\cup\{l\}})$ (contraction with $a$) satisfies $\partial h+h\partial=g$; this is the standard homotopy
\cite[\S17]{Eisenbud1995}. (b) This is the depth sensitivity of the Koszul complex \cite[\S1.6]{BrunsHerzog1993}
(stated there for homology; $K(r;M)$ in cohomological degree $n$ is the Koszul homology complex in degree $|\kappa|-n$), whose proof (induction on
$d$, using (a) and the exact sequence of Koszul complexes for $M\xrightarrow{f_1}M\to M/f_1M$) applies to arbitrary modules; the development
proves it in this way. (c) The first assertion is clear. $K(r;M)$ is the tensor product of $K(r\circ e;M)$ with the Koszul complex of the remaining
elements; filtering by the number of indices outside the image of $e$, each graded piece is a sum of shifted copies of $K(r\circ e;M)$, exact
below $d$; this gives exactness below $d$ and, by the same filtration (the standard staircase argument for a filtered complex), (d) and (e): in
(d) one extends the component $m$ step by step to a cocycle, the obstructions lying in $Z^{<d}$ of shifted copies of $K(r\circ e;M)$ or, at the
last step, being killed by the hypothesis $r_lm\in(r\circ e)M$; in (e) one removes the component at $e(\{1..d\})$ by a coboundary and then
proceeds by descending induction on the filtration. (f) Since the $r\circ e'$ lie in the ideal generated by the $r\circ e$, the matrix $N$
defines a morphism from the Koszul complex of $r\circ e'$ to that of $r\circ e$ which is $\bigwedge^dN$, i.e.\ multiplication by $\det N$, in top
degree; applying (e) to the cocycle $c$ corrected by this comparison gives the claim. In the development, (f) is deduced from the identity
\lref{UpperTheory.Gal.productCup_of_system.det_of} (the maximal minors of a matrix $\rho$ form a Koszul cocycle for the family
$\sum_b\rho_{lb}f_b$).
\end{proof}

\subsubsection*{Double complexes with exact columns}

\begin{lemma}[Gluing and boundary criterion]\label{pf:hlf-chart:glue}
Let $T$ be a good system with admissible Koszul data $(n,z)$, and assume that every column complex of $\mathrm{Tot}=\mathrm{Tot}_T(n,z)$ is exact
below $d$.
\begin{enumerate}[label=\textup{(\alph*)},nosep]
\item Let $c^0_{ij}$ ($i,j\in[d]$) be $d$-cocycles of the $(\{i\},\{j\})$-column complexes such that $c^0_{ij}-c^0_{i'j}$ is a coboundary of the
  $(\{i,i'\},\{j\})$-column complex and $c^0_{ij}-c^0_{ij'}$ one of the $(\{i\},\{j,j'\})$-column complex. Then there is a cocycle
  $c\in Z^{d+2}(\mathrm{Tot})$ with $\mathrm{col}_{\{i\},\{j\}}(c)=c^0_{ij}$ for all $i,j$ and $c_{I\sqcup\bar J\sqcup L}=0$ for $|L|>d$.
\item Let $c\in Z^N(\mathrm{Tot})$ with $c_{I\sqcup\bar J\sqcup L}=0$ for $|L|>d$, and assume that $\mathrm{col}_{I,J}(c)$ is a coboundary of the
  $(I,J)$-column complex whenever $|I|+|J|=N-d$. Then $c\in B^N(\mathrm{Tot})$.
\end{enumerate}
\formallabel{UpperTheory.Gal.productCup_of_system.coordinateSystem_degree_two_of_universalDerivation}\formallabel{UpperTheory.Gal.productCup_of_system.coordinateSystem_degree_one_of_universalDerivation}
\end{lemma}

\begin{proof}
Filter $\mathrm{Tot}$ by the Koszul degree $|L|$. Since the columns are exact below $d$, the usual staircase argument for double complexes
(\cite[\S5.6]{Weibel1994}; equivalently, the spectral sequence of the filtration degenerates to the \v{C}ech complex of the column cohomologies
in Koszul degree $d$) gives: (a) the cochain $c^0$ (placed in the columns with $|I|=|J|=1$) has \v{C}ech differential a sum of column
coboundaries by hypothesis; choosing column primitives and correcting step by step towards lower Koszul degree produces $c$, each obstruction
being a column cocycle of Koszul degree $<d$, hence a column coboundary; (b) one subtracts from $c$ the total differential of column primitives of
its lowest-\v{C}ech-degree columns, which removes those columns at the cost of terms of lower Koszul degree, and concludes by induction, the
remaining column cocycles being in Koszul degree $<d$ and hence coboundaries. The development proves the corresponding statements for arbitrary
weighted complexes with exact columns (\lref{UpperTheory.Gal.productCup_of_system.localTwo_of_gl2} and \lref{UpperTheory.Gal.productCup_of_system.localOne_of_gl1}, based on the filtration criteria \lref{UpperTheory.Gal.productCup_of_system.filtrationCocycleLiftCriterion_of_wdBasic} and \lref{UpperTheory.Gal.productCup_of_system.filtrationBoundaryCriterion_of_wdBasic})
and specializes them.
\end{proof}

\begin{lemma}[K\"unneth decomposition]\label{pf:hlf-chart:kundec}
Let $(A_s)$, $(B_s)$ be increasing families of subspaces of $E$, $E'$ indexed by the nonempty subsets of $[d]$, and let
$t\in Z^{d+2}(\mathrm{Mix}(A\otimes B))$, where $(A\otimes B)_{I,J}$ is the image of $A_I\otimes B_J$ in $E\otimes_\C E'$. Then there are finitely
many triples $(p_k,a_k,b_k)$ with $p_k\le d$, $a_k\in Z^{d-p_k}(\hlfcC^\bullet(A))$, $b_k\in Z^{p_k}(\hlfcC^\bullet(B))$ such that
\[
  t-\sum_k a_k\boxtimes b_k\in B^{d+2}(\mathrm{Mix}(A\otimes B)),\qquad (a\boxtimes b)_{I\sqcup\bar J}=a_I\otimes b_J .
\]
\formallabel{UpperTheory.Gal.productCup_of_system.kunDec_of}
\end{lemma}

\begin{proof}
$\mathrm{Mix}(A\otimes B)$ is, up to the shift by $2$ and signs, the tensor product over $\C$ of the complexes $\hlfcC^\bullet(A)$ and $\hlfcC^\bullet(B)$.
Over a field, the K\"unneth map $\bigoplus_{p+q=n}\hlfcH^p(A)\otimes\hlfcH^q(B)\to H^n(\hlfcC^\bullet(A)\otimes\hlfcC^\bullet(B))$ is an isomorphism
\cite[Thm.~3.6.3]{Weibel1994}; surjectivity in total degree $d$ gives the decomposition (finitely many terms, as every element of a tensor
product is a finite sum of pure tensors).
\end{proof}

\subsubsection*{Edge maps}
For a cochain $y$ of $\hlfcC^\bullet(M)$ on $[d]$ define the cochains $\mathrm{edge}_1(y)$ and $\mathrm{edge}_2(y)$ of a mixed complex by
$\mathrm{edge}_1(y)_{I\sqcup\bar J}=y_I$ if $|J|=1$ and $0$ otherwise, and $\mathrm{edge}_2(y)_{I\sqcup\bar J}=y_J$ if $|I|=1$ and $0$ otherwise.
For families $a$ on $[d]$ and $b$ on $[d]$ (or on $\bar{[d]}$) and a bilinear $\mu$, $\mathrm{join}_\mu(a,b)_{I\sqcup\bar J}=\mu(a_I,b_J)$.

\begin{lemma}[Edge relations]\label{pf:hlf-chart:edgerel}
Let $\gamma$ be a $p$-cocycle of the \v{C}ech complex of a family $(G_s)$ indexed by the subsets of $[d]\sqcup[d]$, write
$\gamma^{(1)}=(\gamma_{I\sqcup\emptyset})_I$ and $\gamma^{(2)}=(\gamma_{\emptyset\sqcup\bar J})_J$, and let $\mathrm{Mix}(G)$ be the mixed complex of
$G_{I,J}=G_{I\sqcup\bar J}$. Then $\mathrm{edge}_2(\gamma^{(2)})-(-1)^p\,\mathrm{edge}_1(\gamma^{(1)})\in B^{p+2}(\mathrm{Mix}(G))$. In particular, for
a $p$-cocycle $\alpha$ of an increasing family $(G_s)_{s\subseteq[d]}$, $\mathrm{edge}_2(\alpha)-(-1)^p\mathrm{edge}_1(\alpha)$ is a coboundary of
the mixed complex of $(G_{I\cup J})$.
\formallabel{UpperTheory.Gal.productCup_of_system.edgeRel_of}
\end{lemma}

\begin{proof}
Let $\eta$ be the $(p+1)$-cochain of $\mathrm{Mix}(G)$ with $\eta_{I\sqcup\bar J}=\gamma_{I\sqcup\bar J}$ ($I,J\ne\emptyset$). For $I,J\ne\emptyset$
with $|I|+|J|=p+2$, the cocycle condition $(\delta\gamma)_{I\sqcup\bar J}=0$ differs from $(\partial\eta)_{I\sqcup\bar J}$ exactly by the terms
in which the removed index empties $I$ or $J$: removing the only element of $I$ (sign $+1$, as it is the smallest index) gives $\gamma_{\emptyset\sqcup\bar J}$,
and removing the only element of $J$ (sign $(-1)^{|I|}=(-1)^{p+1}$) gives $\gamma_{I\sqcup\emptyset}$. Hence
$\partial\eta=-\mathrm{edge}_2(\gamma^{(2)})+(-1)^p\mathrm{edge}_1(\gamma^{(1)})$. For the second statement apply this to the pull-back
$\gamma$ of $\alpha$ along the map $[d]\sqcup[d]\to[d]$ which is the identity on both copies: $\gamma_{I\sqcup\bar J}=\varepsilon\,\alpha_{I\cup J}$
if $I\cap J=\emptyset$, where $\varepsilon$ is the sign of the permutation putting the elements of $I\sqcup J$ in increasing order, and
$\gamma_{I\sqcup\bar J}=0$ otherwise. Identifying cochains indexed by subsets with alternating cochains, $\gamma$ is the image of $\alpha$
under the refinement map, which commutes with the differentials; so $\gamma$ is a $p$-cocycle (with $\gamma_{I\sqcup\bar J}\in G_{I\cup J}$),
and $\gamma^{(1)}=\gamma^{(2)}=\alpha$.
\end{proof}

\subsection{The diagonal}\label{pf:hlf-chart:ss:diag}

Throughout this subsection we assume \ref{pf:hlf-chart:H1} for $K$ (only $K$ is involved) and fix $a\le d$, a top coefficient $\tau$ and a residue
functional $\rho$.

\subsubsection*{The product system and kernels}
The \emph{product system} $T_Z$ has $S=K\otimes_\C K$ acting on $E=\Omega^\bullet_K\otimes_\C\Omega^\bullet_K$ factorwise, $X_i=x_i\otimes1$,
$Y_i=1\otimes x_i$, $R_{I,J}=B_I\otimes B_J$ and $M_{I,J}=\Omega^{d-a}(B_I)\otimes\Omega^a(B_J)$ (images in $E$; for $I$ or $J$ empty these play no
role). It is good by \cref{pf:hlf-chart:rybasic,pf:hlf-chart:omloc} and since $x_m/x_i\in A_I\subseteq B_I$ for $i\in I$
(\lref{UpperTheory.Gal.productCup_of_system.goodZ_of}). Its mixed complex $\mathrm{Mix}(T_Z)$ is the \v{C}ech complex of the product cover of
$X\times X$ for the sheaf $\Omega^{d-a}\boxtimes\Omega^a$. Let $\hlfcmult\colon K\otimes K\to K$ be the multiplication.

For $m=\sum_k\beta_k\otimes\beta'_k\in E$ and $w\in\Omega^\bullet_K$ put
\[
  \mu_\tau(m,w)=\sum_k\Tr_{K/F}\bigl(\tau(\beta_k\wedge w)\bigr)\otimes\beta'_k\in F\otimes_\C\Omega^\bullet_K,
\]
and let $\hlfcmult_F,\hlfcres_\rho\colon F\otimes_\C\Omega^\bullet_K\to\Omega^\bullet_K$ be $f\otimes\beta\mapsto f\beta$ and $f\otimes\beta\mapsto\rho(f)\beta$. We say that
$m$ \emph{acts as} $\gamma\in K$ if $\hlfcmult_F(\mu_\tau(m,w))=\gamma w$ for all $w\in\Omega^a_K$. Note that $\hlfcres_\rho\mu_\tau(\beta\otimes\beta',w)=\theta(\beta\wedge w)\beta'$.
For a cochain $t$ of $\mathrm{Mix}(T_Z)$ and a cochain $\alpha$ of $\hlfcC^\bullet(B_\bullet,a)$ define the \emph{transform}
\begin{equation}\label{eq:pf:hlf-chart:transform}
  (t\star\alpha)_J=\sum_{i\in[d]}\hlfcres_\rho\,\mu_\tau\bigl(t_{\{0..i\}\sqcup\bar J},\,\alpha_{\{i..d\}}\bigr)\qquad(J\subseteq[d]).
\end{equation}
For cochains $e$ of $\hlfcC^\bullet(B_\bullet,d-a)$ and $f$ of $\hlfcC^\bullet(B_\bullet,a)$ we have $(e\boxtimes f)\star\alpha=\langle e,\alpha\rangle f$
(compare \eqref{eq:pf:hlf-chart:pair}).

For a set $L$ of $d$ ordered pairs $l=(l_1,l_2)\in[d]^2$ let
\[
  \kappa_L=\det\begin{pmatrix}e_0\\ x_{l_1}x_{l_2}\,(e_{l_1}-e_{l_2})\ \ (l\in L)\end{pmatrix}\in F\subseteq K,
\]
the determinant of the $(d+1)\times(d+1)$ matrix whose first row is the unit vector $e_0$ and whose other rows, in increasing order of $l$, are
the rows $x_{l_1}x_{l_2}(e_{l_1}-e_{l_2})$; it is the image under $\hlfcmult$ of the analogous determinant formed with the rows of the derivatives of
the minors $\mu_l=X_{l_1}Y_{l_2}-X_{l_2}Y_{l_1}$ (\lref{UpperTheory.Gal.productCup_of_system.kapInd_of}). For $L=D_0=\{(m,0):1\le m\le d\}$,
$\kappa_{D_0}=x_0^d\prod_{m\ge1}x_m=x_0^{2d}\prod_{l\in[d]}(x_l/x_0)$.

\begin{definition}[Diagonal data]\label{pf:hlf-chart:def:zf}
\emph{Diagonal Koszul data} are admissible Koszul data $(n,z)$ for $T_Z$, indexed by a finite totally ordered set $\kappa$, such that
$\hlfcmult(z_l)=0$ for all $l$, and such that for every $i\in[d]$ the elements $b\in B_{\{i\}}$ for which $z_l=(X_iY_i)^N(b\otimes1-1\otimes b)$ and
$n_l=(N,N)$ for some $l$ and $N$ generate $B_{\{i\}}$ as a $\C$-algebra. A cocycle $c$ of $\mathrm{Tot}_{T_Z}(\mu)$, the \v{C}ech--Koszul complex of
the minors, is \emph{diagonal} if $c_{\{i\}\sqcup\{\bar j\}\sqcup L}$ acts as $(-1)^d\kappa_L$ for all $i,j\in[d]$ and all sets $L$ of $d$ pairs.
\formal{UpperTheory.Gal.productCup_of_system.ZF}\formal{UpperTheory.Gal.productCup_of_system.IsDiag}
\end{definition}

Geometrically, the minors cut out the diagonal of $\mathbb P^d\times\mathbb P^d$, hence $X\times_{\mathbb P^d}X$ in $X\times X$, and the $z_l$ cut out the
diagonal of $X$ inside $X\times_{\mathbb P^d}X$. A diagonal cocycle is a \v{C}ech--Koszul representative of the class of the diagonal, normalized
at the local level by the condition that its top Koszul components act as the identity (up to the factor $(-1)^d\kappa_L$).

\begin{lemma}[Diagonal data exist]\label{pf:hlf-chart:zfam}
There are diagonal Koszul data $(n,z)$ indexed by $\kappa=\{1,\dots,M\}$.
\formallabel{UpperTheory.Gal.productCup_of_system.zFam_of}
\end{lemma}

\begin{proof}
For each $i$ choose finitely many $\C$-algebra generators $b$ of $B_{\{i\}}$ (\cref{pf:hlf-chart:ryft}). For each such $b$ and each $i'\in[d]$,
$b\in B_{\{i\}}\subseteq B_{\{i,i'\}}$, so $b\,(x_i/x_{i'})^{N}\in B_{\{i'\}}$ for large $N$ (\cref{pf:hlf-chart:rybasic}(c)); choose $N=N(b)$ working for all
$i'$ and put $z=(X_iY_i)^N(b\otimes1-1\otimes b)$, $n=(N,N)$. Then $\hlfcmult(z)=x_i^{2N}(b-b)=0$, and for all $i',j'$
\[
  \frac{z}{X_{i'}^NY_{j'}^N}=\Bigl(\frac{x_i}{x_{i'}}\Bigr)^Nb\otimes\Bigl(\frac{x_i}{x_{j'}}\Bigr)^N-\Bigl(\frac{x_i}{x_{i'}}\Bigr)^N\otimes\Bigl(\frac{x_i}{x_{j'}}\Bigr)^Nb
  \in B_{\{i'\}}\otimes B_{\{j'\}},
\]
so the data are admissible. Enumerate all these elements.
\end{proof}

\begin{lemma}[A diagonal cocycle]\label{pf:hlf-chart:zex}
Let $(n,z)$ be diagonal Koszul data. There is a cocycle $c\in Z^{d+2}(\mathrm{Tot}_{T_Z}(\mu\sqcup z))$, for the Koszul data consisting of the minors
(bidegree $(1,1)$) followed by $(n,z)$, whose restriction to the minors is diagonal.
\formallabel{UpperTheory.Gal.productCup_of_system.zEx_of}
\end{lemma}

\begin{proof}
We use the following results of the development.
(i) For $i\in I$, $j\in J$, the Koszul complex of the $d$ dehomogenized minors $r_{(m,j)}(i,j)$, $m\ne j$, on $M_{I,J}$ is exact below $d$,
and an element of $M_{I,J}$ which acts as $0$ lies in the submodule generated by these minors
(\lref{UpperTheory.Gal.productCup_of_system.colZ_of}); in the chart $X_iY_j\ne0$ the elements $r_{(m,j)}(i,j)$ are, up to units, the
differences $s_m\otimes1-1\otimes s_m$ of the affine coordinates, which form a regular sequence on $M_{I,J}$, a lattice over $A_I\otimes A_J$
(\lref{UpperTheory.Gal.productCup_of_system.latOm_of} and \lref{UpperTheory.Gal.productCup_of_system.regLat_of}), and the second statement
is the trace duality of \cref{pf:hlf-chart:dual} (\lref{UpperTheory.Gal.productCup_of_system.dualForm_of}).
(ii) For all $i,j$ there is a $d$-cocycle $c^0_{ij}$ of the $(\{i\},\{j\})$-column complex of the Koszul data $\mu\sqcup z$ whose components at the
sets $L$ of $d$ minors act as $(-1)^d\kappa_L$ (\lref{UpperTheory.Gal.productCup_of_system.zCol_of}): it is obtained from the Koszul cocycle
formed by the maximal minors of the matrix of derivatives of the $\mu_l$ (\lref{UpperTheory.Gal.productCup_of_system.thMu_of}), multiplied
by a kernel which acts as $\hlfcmult$ of the relevant determinant (\lref{UpperTheory.Gal.productCup_of_system.idIJ_of}), and extended to the
$z_l$ by \cref{pf:hlf-chart:koszul}(d) (the products of $z_l$ with such a kernel act as $\hlfcmult(z_l)\cdot(\ )=0$, the kernels compatible with the
action forming a module over $K\otimes K$, \lref{UpperTheory.Gal.productCup_of_system.actsLin_of}).

By (i) and \cref{pf:hlf-chart:sysbasic}(e), \cref{pf:hlf-chart:koszul}(c), all column complexes of $\mathrm{Tot}_{T_Z}(\mu\sqcup z)$ are exact below $d$.
On an overlap $(\{i,i'\},\{j\})$ or $(\{i\},\{j,j'\})$ consider the difference of the two cocycles $c^0$ (regarded in the column of the overlap,
\cref{pf:hlf-chart:sysbasic}(e)). Its component at $D_j$ is the difference of two elements acting as $(-1)^d\kappa_{D_j}$, hence acts as $0$, hence
lies in the submodule generated by the minors indexed by $D_j$ by (i); so the difference is a column coboundary by \cref{pf:hlf-chart:koszul}(e).
\Cref{pf:hlf-chart:glue}(a) glues the $c^0_{ij}$ to a cocycle $c$ of degree $d+2$ with $(\{i\},\{\bar j\})$-columns $c^0_{ij}$. Its restriction to the
minors (\cref{pf:hlf-chart:sysbasic}(d)) has components $c_{\{i\}\sqcup\{\bar j\}\sqcup L}=(c^0_{ij})_L$ for sets $L$ of $d$ minors, so it is diagonal.
\end{proof}

\begin{lemma}[The diagonal acts as the identity]\label{pf:hlf-chart:zact}
Let $c$ be a diagonal cocycle of $\mathrm{Tot}_{T_Z}(\mu)$ of degree $d+2$. Then for every $q$ and every $\alpha\in Z^q(\hlfcC^\bullet(B_\bullet,a))$,
$\alpha-\mathrm{top}(c)\star\alpha\in B^q(\hlfcC^\bullet(B_\bullet,a))$.
\formallabel{UpperTheory.Gal.productCup_of_system.zAct_of}
\end{lemma}

\begin{proof}
The proof compares $c$ with an explicit cocycle over $\mathbb P^d\times X$. Let $T_P$ be the system with $S=F\otimes K$ acting on $F\otimes\Omega^\bullet_K$,
$X_i=x_i\otimes1$, $Y_i=1\otimes x_i$, $R_{I,J}=A_I\otimes B_J$ and $M_{I,J}=\Omega_F(I)\otimes\Omega^a(B_J)$. Let $\alpha\in Z^q$. The development
proves:
(i) the cochain $\Theta$ with
$\Theta_{I\sqcup\bar J\sqcup L}=(-1)^{q(|J|+|L|)}\sum_{i\in I}\mu_\tau(c_{\{i'\in I:i'\le i\}\sqcup\bar J\sqcup L},\alpha_{\{i'\in I:i'\ge i\}})$
is a cocycle of $\mathrm{Tot}_{T_P}(\mu)$ of degree $d+q+2$, vanishing in Koszul degrees $>d$, and $\hlfcmult_F(\Theta_{b\sqcup L})=
\varepsilon\,\kappa_L\,\mathrm{edge}_1(\alpha)_b$ for $|L|=d$, with $\varepsilon=(-1)^{q(d+1)+d}$ (\lref{UpperTheory.Gal.productCup_of_system.pCup_of};
this uses that $c$ is diagonal);
(ii) there is a cocycle $\Theta'$ of $\mathrm{Tot}_{T_P}(\mu)$ of the same degree, vanishing in Koszul degrees $>d$, with
$\hlfcmult_F(\Theta'_{b\sqcup L})=\varepsilon'\,\kappa_L\,\mathrm{edge}_2(\alpha)_b$ for $|L|=d$, $\varepsilon'=(-1)^{(q+1)d}$, and
$\Theta'_{[d]\sqcup\bar J\sqcup\emptyset}=1\otimes\alpha_J$ for all $J$ (\lref{UpperTheory.Gal.productCup_of_system.pTheta_of}); it is built from
the Koszul cocycle of maximal minors over $\mathbb P^d$, which is explicit (\lref{UpperTheory.Gal.productCup_of_system.mJ_of});
(iii) a cocycle of $\mathrm{Tot}_{T_P}(\mu)$ vanishing in Koszul degrees $>d$ whose top Koszul components are $\kappa_L$ times a coboundary of the
mixed complex of $(\Omega^a(B_{I\cup J}))$ is a coboundary (\lref{UpperTheory.Gal.productCup_of_system.pCmp_of}; this is
\cref{pf:hlf-chart:glue}(b) combined with the exactness of the columns, \lref{UpperTheory.Gal.productCup_of_system.colP_of});
(iv) if $t$ is a coboundary of $\mathrm{Mix}(T_P)$ then $(\hlfcres_\rho t_{[d]\sqcup\bar J})_J$ is a coboundary of $\hlfcC^\bullet(B_\bullet,a)$
(\lref{UpperTheory.Gal.productCup_of_system.pInt_of}; by \cref{pf:hlf-chart:rho}(b), $\rho$ kills the components over index sets $I\ne[d]$).

By \cref{pf:hlf-chart:edgerel}, $\mathrm{edge}_2(\alpha)-(-1)^q\mathrm{edge}_1(\alpha)$ is a coboundary of the mixed complex. Since
$\varepsilon\varepsilon'=(-1)^q$, the top Koszul components of $\Theta-\Theta'$ satisfy
$\hlfcmult_F((\Theta-\Theta')_{b\sqcup L})=-\varepsilon'\kappa_L(\mathrm{edge}_2(\alpha)-(-1)^q\mathrm{edge}_1(\alpha))_b$, so (iii) shows that
$\Theta-\Theta'$ is a coboundary of $\mathrm{Tot}_{T_P}(\mu)$. Its top part is then a coboundary of $\mathrm{Mix}(T_P)$ (\cref{pf:hlf-chart:sysbasic}(b)),
and by (iv) $(\hlfcres_\rho(\Theta-\Theta')_{[d]\sqcup\bar J\sqcup\emptyset})_J$ is a coboundary of $\hlfcC^\bullet(B_\bullet,a)$. Since $\rho(1)=1$, the
second term is $\alpha$. In the first, only $|J|=q+1$ contributes (the components of $c$ have $i+1+|J|=d+2$ and those of $\alpha$ have $d-i+1=q+1$),
the sign is $(-1)^{q(q+1)}=1$, and the term is $\mathrm{top}(c)\star\alpha$ by \eqref{eq:pf:hlf-chart:transform}.
\end{proof}

\begin{lemma}[Transforms of coboundaries]\label{pf:hlf-chart:rcuptop}
If $t\in B^{d+2}(\mathrm{Mix}(T_Z))$ and $\alpha\in Z^q(\hlfcC^\bullet(B_\bullet,a))$, then $t\star\alpha\in B^q(\hlfcC^\bullet(B_\bullet,a))$.
\formallabel{UpperTheory.Gal.productCup_of_system.rCupTop_of}
\end{lemma}

\begin{proof}
Let $t=\partial s$. Expanding $\partial s$ by the decomposition of the differential along $[d]\sqcup\bar{[d]}$ (\cref{pf:hlf-chart:def:wcx}), the
part of $\partial s$ differentiating in the second copy gives, after the transform, $\pm\delta(s\star\alpha)$; the part differentiating in the first
copy gives, by the computation of \cref{pf:hlf-chart:adj} with $J$ as a passive parameter and $\delta\alpha=0$, a sum of terms
$\hlfcres_\rho\mu_\tau(s_{I'\sqcup\bar J},\alpha_{I''})$ with $I'\cup I''\subseteq[d]\smallsetminus\{i\}$ for some $i$. Such terms vanish: their first
factors are $\Tr_{K/F}\tau(\beta\wedge w)$ with $\beta\wedge w\in\Omega^d(B_{[d]\smallsetminus\{i\}})$, which lie in $\Omega_F([d]\smallsetminus\{i\})$
(\cref{pf:hlf-chart:trreg}(a)) and are killed by $\rho$. Hence $t\star\alpha=\pm\delta(s\star\alpha)$.
\end{proof}

\begin{lemma}[Dual decomposition of the diagonal]\label{pf:hlf-chart:zid}
Let $c$ be a diagonal cocycle of $\mathrm{Tot}_{T_Z}(\mu)$ of degree $d+2$. There are finitely many triples $(p_k,e_k,f_k)$ with $p_k\le d$,
$e_k\in Z^{d-p_k}(\hlfcC^\bullet(B_\bullet,d-a))$ and $f_k\in Z^{p_k}(\hlfcC^\bullet(B_\bullet,a))$ such that
\begin{equation}\label{eq:pf:hlf-chart:topdec}
  \mathrm{top}(c)-\sum_ke_k\boxtimes f_k\in B^{d+2}(\mathrm{Mix}(T_Z)),
\end{equation}
and every $\alpha\in Z^q(\hlfcC^\bullet(B_\bullet,a))$ satisfies \eqref{eq:pf:hlf-chart:dualdec}.
\formallabel{UpperTheory.Gal.productCup_of_system.zId_of}
\end{lemma}

\begin{proof}
$\mathrm{top}(c)$ is a cocycle of $\mathrm{Mix}(T_Z)=\mathrm{Mix}(A\otimes B)$ with $A_I=\Omega^{d-a}(B_I)$, $B_J=\Omega^a(B_J)$ (\cref{pf:hlf-chart:sysbasic}(b)),
and these families are increasing (\cref{pf:hlf-chart:rybasic,pf:hlf-chart:omloc}). \Cref{pf:hlf-chart:kundec} gives the triples and
\eqref{eq:pf:hlf-chart:topdec}. Let $\alpha\in Z^q$. By \cref{pf:hlf-chart:zact}, $\alpha\equiv\mathrm{top}(c)\star\alpha$; by
\cref{pf:hlf-chart:rcuptop} applied to the coboundary in \eqref{eq:pf:hlf-chart:topdec}, $\mathrm{top}(c)\star\alpha\equiv\sum_k(e_k\boxtimes f_k)\star\alpha
=\sum_k\langle e_k,\alpha\rangle f_k$, modulo $B^q$. Finally $\langle e_k,\alpha\rangle=0$ unless $p_k=q$: the summand $j$ of
\eqref{eq:pf:hlf-chart:pair} needs $j+1=d-p_k+1$ and $d-j+1=q+1$.
\end{proof}

\subsection{The correspondence and the localization}\label{pf:hlf-chart:ss:graph}

In this subsection we assume \ref{pf:hlf-chart:H1}--\ref{pf:hlf-chart:H3}.

\subsubsection*{The graph system}
The \emph{graph system} $T_Y$ has $S=K'$ acting on $E=\Omega^\bullet_{K'}$ by multiplication, $X_i=x_i$, $Y_i=j_2(x_i)$, $R_{I,J}=B'_{I,J}$ and
$M_{I,J}=\Omega^d(B'_{I,J})$; it is good by \cref{pf:hlf-chart:rybasic,pf:hlf-chart:omloc}, as $x_m/x_i\in B'_I$ and $j_2(x_m/x_j)\in B''_J$ for
$i\in I$, $j\in J$. Its mixed complex is the \v{C}ech complex of $\Omega^d_Y$ for the product of the covers $(c_1^{-1}U_i)$ and $(c_2^{-1}U_j)$. The
\emph{graph maps} are the $\C$-algebra map $g\colon K\otimes K\to K'$, $b\otimes b'\mapsto j_1(b)j_2(b')$, and the $\C$-linear map
\[
  \gamma\colon\Omega^\bullet_K\otimes\Omega^\bullet_K\to\Omega^\bullet_{K'},\qquad\beta\otimes\beta'\mapsto j_1^*\beta\wedge j_2^*\beta'.
\]
We extend $\gamma$ to cochains componentwise.

\begin{lemma}[Pull-back along the graph]\label{pf:hlf-chart:gmap}
For $a\le d$: $g$ maps the minors of $T_Z$ to the minors $x_mj_2(x_{m'})-x_{m'}j_2(x_m)$ of $T_Y$; $\gamma(s\cdot m)=g(s)\gamma(m)$;
$g(B_I\otimes B_J)\subseteq B'_{I,J}$; $\gamma$ maps $\mathrm{Mix}(T_Z)$-coboundaries to $\mathrm{Mix}(T_Y)$-coboundaries; and for all Koszul data $(n,z)$ of
$T_Z$, $\gamma$ maps cocycles of $\mathrm{Tot}_{T_Z}(n,z)$ to cocycles of $\mathrm{Tot}_{T_Y}(n,g(z))$.
\formallabel{UpperTheory.Gal.productCup_of_system.gMap_of}
\end{lemma}

\begin{proof}
The first three statements follow from the definitions and \cref{pf:hlf-chart:rybasic}(b), since $j_1^*$ and $j_2^*$ are ring homomorphisms. In
particular $\gamma(\Omega^{d-a}(B_I)\otimes\Omega^a(B_J))\subseteq\Omega^{d-a}(B'_{I,J})\wedge\Omega^a(B'_{I,J})\subseteq\Omega^d(B'_{I,J})$, and $\gamma$
maps the modules $u_L\,M^Z_{I,J}$ to $g(u_L)\,M^Y_{I,J}$, with $g(X_i)=x_i$ and $g(Y_i)=j_2(x_i)$. As $\gamma$ is linear, applied componentwise, and
satisfies $\gamma(z_l m)=g(z_l)\gamma(m)$, it commutes with the weighted differentials.
\end{proof}

\begin{lemma}[Cup products and edges]\label{pf:hlf-chart:lcup}
Let $\mu\colon E\times E'\to E''$ be bilinear, $(A_I)$ a family of subspaces of $E$ indexed by the subsets of $[d]$, $(G_s)$ one of $E'$ indexed by
the subsets of $[d]\sqcup[d]$, and $(H_{I,J})$ one of $E''$, such that $\mu(A_{I'},G_s)\subseteq H_{I,J}$ whenever $\emptyset\ne I'\subseteq I$,
$\emptyset\ne s\subseteq I\sqcup\bar J$ and $J\ne\emptyset$. Let $a\in Z^q(\hlfcC^\bullet(A))$ and $\gamma\in Z^p(\hlfcC^\bullet(G))$, and put
$\gamma^{(1)}=(\gamma_{I\sqcup\emptyset})_I$, $\gamma^{(2)}=(\gamma_{\emptyset\sqcup\bar J})_J$. Then $a\cup_\mu\gamma^{(1)}$ is a $(p+q)$-cocycle and
\[
  \mathrm{join}_\mu(a,\gamma^{(2)})-(-1)^p\,\mathrm{edge}_1(a\cup_\mu\gamma^{(1)})\in B^{p+q+2}(\mathrm{Mix}(H)).
\]
\formallabel{UpperTheory.Gal.productCup_of_system.lCupRel_of}
\end{lemma}

\begin{proof}
$\delta(a\cup\gamma^{(1)})=\delta a\cup\gamma^{(1)}\pm a\cup\delta\gamma^{(1)}=0$ by the Leibniz rule. Let $\tilde a$ be the cochain on $[d]\sqcup[d]$
with $\tilde a_{I\sqcup\emptyset}=a_I$ and $\tilde a_{I\sqcup\bar J}=0$ for $J\ne\emptyset$. Then $(\delta\tilde a)_{I\sqcup\{\bar j\}}=(-1)^{|I|}a_I$ and
all other components of $\delta\tilde a$ vanish, so $\Gamma=\tilde a\cup_\mu\gamma$ satisfies
$(\delta\Gamma)_{I\sqcup\bar J}=(\delta\tilde a\cup\gamma)_{I\sqcup\bar J}=(-1)^{|I|}\mu(a_I,\gamma_{\emptyset\sqcup\bar J})$ for $J\ne\emptyset$ (the only
splitting point is $\min\bar J$), while $\Gamma_{\emptyset\sqcup\bar J}=0$ and $\Gamma_{I\sqcup\emptyset}=(a\cup\gamma^{(1)})_I$. Let $\eta$ be the
cochain of $\mathrm{Mix}(H)$ with $\eta_{I\sqcup\bar J}=\Gamma_{I\sqcup\bar J}=\sum_{i\in I}\mu(a_{\{i'\in I:i'\le i\}},\gamma_{\{i'\in I:i'\ge i\}\sqcup\bar J})$
for $I,J\ne\emptyset$; it lies in $\mathrm{Mix}(H)$ by the hypothesis on $\mu$. As in the proof of \cref{pf:hlf-chart:edgerel},
$(\partial\eta)_{I\sqcup\bar J}$ is $(\delta\Gamma)_{I\sqcup\bar J}$ minus the terms in which the removed index empties $I$ or $J$; with $|I|=q+1$ in the
first and $|I|=p+q+1$ in the last term this gives
$\partial\eta=(-1)^{q+1}\bigl(\mathrm{join}_\mu(a,\gamma^{(2)})-(-1)^p\mathrm{edge}_1(a\cup\gamma^{(1)})\bigr)$.
\end{proof}

\begin{lemma}[Exact rows]\label{pf:hlf-chart:rowex}
Assume \ref{pf:hlf-chart:H1}, \ref{pf:hlf-chart:H2}, and let $\emptyset\ne I\subseteq[d]$, $q\ge0$. The weighted \v{C}ech complex on $[d]$ with modules
$\Omega^q(B'_{I,J})$ for \emph{all} $J\subseteq[d]$ (including $J=\emptyset$, where the module is $\Omega^q(B'_I)$) and weights $1$ is exact.
\formallabel{UpperTheory.Gal.productCup_of_system.rowEx_of}
\end{lemma}

\begin{proof}
All modules are $B'_I$-modules (\cref{pf:hlf-chart:omloc}) and $B'_I$ is a finitely generated $\C$-algebra (\cref{pf:hlf-chart:ryft}), so it suffices to
prove exactness after localization at the maximal ideal $\ker e$ of every character $e$ of $B'_I$
(\cite[Prop.~3.9]{AtiyahMacdonald1969} and the Nullstellensatz; the development proves this local--global principle directly,
\lref{UpperTheory.Gal.productCup_of_system.lG_of}). By \cref{pf:hlf-chart:omloc} the localized modules are $\Omega^q$ of the rings
$(B'_I)_e\cdot B''_J$. Choose $i\in I$; by \ref{pf:hlf-chart:H2} applied to the restriction of $e$ to $B'_{\{i\}}$ there is $j\in[d]$ with
$B''_{\{j\}}\subseteq(B'_{\{i\}})_e\subseteq(B'_I)_e$, and then $(B'_I)_e\cdot B''_{J\cup\{j\}}=(B'_I)_e\cdot B''_J$ for all $J$, since
$B''_{J\cup\{j\}}=B''_J[j_2(x_{j'}/x_j)]$ for $j'\in J$ (\cref{pf:hlf-chart:rybasic}(c)) and $j_2(x_{j'}/x_j)\in B''_{\{j\}}$ (for $J=\emptyset$ the
claim is $B''_{\{j\}}\subseteq(B'_I)_e$). For such a family the complex is contractible: the map $h(c)_J=\pm c_{J\cup\{j\}}$ ($j\notin J$),
$h(c)_J=0$ ($j\in J$), with the sign making $j$ the first index, satisfies $\delta h+h\delta=\mathrm{id}$, as in the proof of
\cite[III, Lemma~4.2]{Hartshorne1977}.
\end{proof}

\begin{lemma}[Edge injectivity]\label{pf:hlf-chart:edgeb}
Assume \ref{pf:hlf-chart:H1}, \ref{pf:hlf-chart:H2}. If $y\in Z^d(\hlfcC^\bullet(B'_\bullet,d))$ and $\mathrm{edge}_1(y)\in B^{d+2}(\mathrm{Mix}(T_Y))$, then
$y\in B^d(\hlfcC^\bullet(B'_\bullet,d))$.
\formallabel{UpperTheory.Gal.productCup_of_system.edgeB_of}
\end{lemma}

\begin{proof}
Let $\mathcal D$ be the weighted \v{C}ech complex on $[d]\sqcup[d]$ with modules $\Omega^d(B'_{I,J})$ for $I\ne\emptyset$ (and all $J$, including $\emptyset$)
and $0$ for $I=\emptyset$, and weights $1$. Filtering by $|I|$, its graded pieces are the complexes of \cref{pf:hlf-chart:rowex}, which are exact; hence
$\mathcal D$ is exact (\cite[Lemma~2.7.3]{Weibel1994}, the filtration being finite). Let $\tilde y$ be the cochain of $\mathcal D$ with
$\tilde y_{I\sqcup\emptyset}=y_I$ and all other components $0$. Since $\delta y=0$, $\partial\tilde y$ has only the components
$(\partial\tilde y)_{I\sqcup\{\bar j\}}=(-1)^{|I|}y_I$, i.e.\ $\partial\tilde y=(-1)^{d+1}\mathrm{edge}_1(y)$, regarded in $\mathcal D$. If
$\mathrm{edge}_1(y)=\partial\xi$ with $\xi$ a cochain of $\mathrm{Mix}(T_Y)\subseteq\mathcal D$, then $\tilde y-(-1)^{d+1}\xi$ is a cocycle of $\mathcal D$, hence
$=\partial\zeta$. Taking components with $J=\emptyset$ (where $\xi$ vanishes and only differentials in the first copy occur) gives
$y=\delta\bigl((\zeta_{I\sqcup\emptyset})_I\bigr)$.
\end{proof}

\subsubsection*{The fibre over the origin}

\begin{lemma}[Points of the fibre]\label{pf:hlf-chart:crt}
Let $x$ define a finite cover of $L$ with $B_{\{0\}}(L)/(t_1,\dots,t_d)$ reduced. Then $\hlfcfib_L$ is finite, and for every function
$v\colon\hlfcfib_L\to\C$ there is $h\in B_{\{0\}}(L)$ with $e(h)=v(e)$ for all $e\in\hlfcfib_L$.
\formallabel{UpperTheory.Gal.productCup_of_system.cRT_of}
\end{lemma}

\begin{proof}
$B=B_{\{0\}}(L)$ is finite over $A_{\{0\}}$ (\cref{pf:hlf-chart:ryft}), so $B/(t)$ is a finite-dimensional reduced $\C$-algebra, hence isomorphic to
$\C^{\hlfcfib_L}$ by the Chinese remainder theorem, its characters being the elements of $\hlfcfib_L$. Lift the element $v$ along the surjection
$B\to B/(t)$.
\end{proof}

Let $dt_\star=dt_1\wedge\dots\wedge dt_d=\bigl(\prod_{l\ge1}t_l\bigr)\omega'_0\in\Omega^d(B'_{\{0\}})$ (in $K'$; similarly in $K$). The \emph{origin data} of
$T_Y$ are the Koszul data $(x_1,\dots,x_d)$ of bidegree $(1,0)$; their dehomogenizations at $(0,j)$ are $t_1,\dots,t_d$, which cut out
$\hlfcfib$ in the chart $B'_{\{0\}}$. For diagonal Koszul data $(n,z)$ of $T_Z$ (\cref{pf:hlf-chart:def:zf}) the \emph{fixed-point data} of $T_Y$ are the
images under $g$ of the minors followed by $(n,g(z))$; they cut out the points where $c_1$ and $c_2$ agree.

\begin{lemma}[Value of the pulled-back diagonal kernel]\label{pf:hlf-chart:valg}
Assume \ref{pf:hlf-chart:H1}--\ref{pf:hlf-chart:H6}, \ref{pf:hlf-chart:H9} and \eqref{eq:pf:hlf-chart:tan}, and let $a\le d$. If
$m\in\Omega^{d-a}(B_{\{0\}})\otimes\Omega^a(B_{\{0\}})$ acts as $\prod_{l\ge1}t_l$, then $\gamma(m)=s\,dt_\star$ for some $s\in K'$ with
$s(e)=[X^a]\det(1+XA_e)$ for every $e\in\mathcal F$ (value in $B'_{\{0\}}$).
\formallabel{UpperTheory.Gal.productCup_of_system.valG_of}
\end{lemma}

\begin{proof}
Write $m=\sum_kw'_k\otimes w_k$. Let $e\in\mathcal F$. By the development: (i) $e$ lies over a point $z$ of the fibre of $B_{\{0\}}$ such that $j_1(v)$
and $j_2(v)$ have the value $v(z)$ at $e$ whenever $v$ has a value at $z$
(\lref{UpperTheory.Gal.productCup_of_system.fixZ_of}, from \ref{pf:hlf-chart:H5} and the definition of fixed points); (ii) the local rings at
points of the fibres are regular with parameters $t_1,\dots,t_d$, and the regular $q$-forms near such a point are the combinations of the
$dt_S=\bigwedge_{l\in S}dt_l$ ($|S|=q$) with coefficients in the local ring
(\lref{UpperTheory.Gal.productCup_of_system.locForms_of} and \lref{UpperTheory.Gal.productCup_of_system.dtBasic_of}); (iii) writing
$w'_k\wedge dt_S=g_{kS}\,dt_\star$ and $w_k=\sum_{|T|=a}b_{kT}\,dt_T$, the values $\alpha_{kS}=g_{kS}(z)$, $\beta_{kT}=b_{kT}(z)$ satisfy
$\sum_k\alpha_{kS}\beta_{kT}=\delta_{ST}$ for $|S|=|T|=a$, because $m$ acts as $\prod t_l$ and $\omega_0=dt_\star/\prod t_l$
(\lref{UpperTheory.Gal.productCup_of_system.idVal_of}, which uses \cref{pf:hlf-chart:rt1} and the trace duality); (iv)
$j_1^*dt_T=dt_T$ and $j_2^*dt_T=\sum_{|S|=|T|}c_{TS}\,dt_S$ with $c_{TS}(e)=\det\bigl((A_e)_{T,S}\bigr)$, by \eqref{eq:pf:hlf-chart:tan}
(\lref{UpperTheory.Gal.productCup_of_system.dtPull_of}). Then, since $j_1^*dt_\star=dt_\star$,
$\gamma(m)=\sum_kj_1^*w'_k\wedge j_2^*w_k=\sum_{k,T,S}j_2(b_{kT})\,c_{TS}\,j_1^*w'_k\wedge dt_S=s\,dt_\star$ with
$s=\sum_{k,T,S}j_1(g_{kS})\,j_2(b_{kT})\,c_{TS}$, all factors having values at $e$; by (i), (iii), (iv),
$s(e)=\sum_{S,T}\bigl(\sum_k\alpha_{kS}\beta_{kT}\bigr)\det((A_e)_{T,S})=\sum_{|S|=a}\det((A_e)_{S,S})=[X^a]\det(1+XA_e)$ (\cref{pf:hlf-chart:minorsum}).
\end{proof}

The following three statements are the local heart of the localization; they are proved in the development by the lemmas of
\cref{pf:hlf-chart:ss:cx} applied to the Koszul complexes of the fixed-point and origin data on $Y$, whose exactness properties
(\lref{UpperTheory.Gal.productCup_of_system.yCols_of} and \lref{UpperTheory.Gal.productCup_of_system.yKE_of}) come from the following facts:
at a character of $B'_{I,J}$ at which all fixed-point data vanish, $J=\{0\}$ (by \ref{pf:hlf-chart:H4}) and the character lies over a point of
$\mathcal F$ (by \ref{pf:hlf-chart:H5}, \ref{pf:hlf-chart:H6}; \lref{UpperTheory.Gal.productCup_of_system.ptG_of}); at $e\in\mathcal F$ the
dehomogenized minors $r_{(a,0)}(0,0)=t_a-j_2(t_a)$ satisfy $t_a-j_2(t_a)=\sum_lN_{al}t_l$ with $N=(\delta_{al}-g_{al})$, a matrix over the local ring
whose determinant has the value $\det(1-A_e)\ne0$, so that conversely the $t_l$ lie locally in the ideal of the $r_{(a,0)}(0,0)$, by
\eqref{eq:pf:hlf-chart:tan}, \ref{pf:hlf-chart:H9} (\lref{UpperTheory.Gal.productCup_of_system.fixLoc_of}); and $t_1,\dots,t_d$ is a regular sequence on the local forms at
points of the fibre, which is reduced by \ref{pf:hlf-chart:H3} (\lref{UpperTheory.Gal.productCup_of_system.locReg_of} and \lref{UpperTheory.Gal.productCup_of_system.locMax_of}).
Exactness and membership in ideals are then checked locally (\lref{UpperTheory.Gal.productCup_of_system.lG_of}).

\begin{lemma}[Local cocycles on $Y$]\label{pf:hlf-chart:ybig}
Assume \ref{pf:hlf-chart:H1}--\ref{pf:hlf-chart:H6}, \ref{pf:hlf-chart:H9} and \eqref{eq:pf:hlf-chart:tan}; let $(n,z)$ be diagonal Koszul data and
$h\in B'_{\{0\}}$.
\begin{enumerate}[label=\textup{(\alph*)},nosep]
\item If $e(h)=0$ for $e\in\hlfcfib\smallsetminus\mathcal F$, there is a cocycle $c_1$ of degree $d+2$ of $\mathrm{Tot}_{T_Y}$ for the fixed-point data
followed by the origin data, whose component at $\{0\}\sqcup\{\bar0\}\sqcup\{\text{all origin indices}\}$ is $(-1)^dx_0^dh\,dt_\star$ and whose
$(\{i\},\{\bar j\})$-columns vanish for $(i,j)\ne(0,0)$ \textup{(\lref{UpperTheory.Gal.productCup_of_system.yBig_of})}.
\item If $e(h)=0$ for $e\in\hlfcfib\smallsetminus\mathcal F$ and $c$ is a cocycle of degree $d+2$ of $\mathrm{Tot}_{T_Y}$ for the origin data, with component
$(-1)^dx_0^dh\,dt_\star$ at $\{0\}\sqcup\{\bar0\}\sqcup\{1..d\}$ and vanishing $(\{i\},\{\bar j\})$-columns for $(i,j)\ne(0,0)$, then
$\mathrm{top}(c)-\mathrm{edge}_1(h\omega'_0)\in B^{d+2}(\mathrm{Mix}(T_Y))$, where $h\omega'_0$ denotes the $d$-cochain of $\hlfcC^\bullet(B'_\bullet,d)$ with
component $h\omega'_0$ at $[d]$ \textup{(\lref{UpperTheory.Gal.productCup_of_system.yCmpT_of})}.
\item If $s\in K'$ has the value $e(h)\det(1-A_e)$ at every $e\in\mathcal F$, $c_1$ is as in \textup{(a)}, and $c'$ is a cocycle of degree $d+2$ of
$\mathrm{Tot}_{T_Y}$ for the fixed-point data with component $(-1)^d(x_0j_2(x_0))^ds\,dt_\star$ at $\{0\}\sqcup\{\bar0\}\sqcup D_0$, then
$\mathrm{top}(c')-\mathrm{top}(c_1)\in B^{d+2}(\mathrm{Mix}(T_Y))$ \textup{(\lref{UpperTheory.Gal.productCup_of_system.yCmpG_of})}.
\end{enumerate}
\end{lemma}

The idea is the transformation law of local residues: (a) and (b) say that the class of the Koszul cocycle
$h\,dt_\star/(t_1\cdots t_d)$ supported on the fibre is the \v{C}ech class $h\omega'_0$ on $\bigcap_iB'_{\{i\}}$ (the residue map in \v{C}ech form), while
(c) changes the regular sequence $t$ into the sequence $t-j_2(t)$ cutting out the fixed points, which multiplies the local component by
$\det(N)$, whose value at $e$ is $\det(1-A_e)$ (\cref{pf:hlf-chart:koszul}(f)). We do not reproduce the bookkeeping of the columns.

\begin{lemma}[Localization]\label{pf:hlf-chart:ymain}
Assume \ref{pf:hlf-chart:H1}--\ref{pf:hlf-chart:H6}, \ref{pf:hlf-chart:H9} and \eqref{eq:pf:hlf-chart:tan}; let $a\le d$, $(n,z)$ diagonal Koszul data
and $c\in Z^{d+2}(\mathrm{Tot}_{T_Z}(\mu\sqcup z))$ a cocycle whose restriction to the minors is diagonal. Then there is $h\in B'_{\{0\}}$ such that
\begin{enumerate}[label=\textup{(\alph*)},nosep]
\item $\gamma(\mathrm{top}(c))-\mathrm{edge}_1(h\omega'_0)\in B^{d+2}(\mathrm{Mix}(T_Y))$;
\item the cochain $h\omega'_0$ (component at $[d]$ only) is a $d$-cocycle of $\hlfcC^\bullet(B'_\bullet,d)$;
\item $\sum_{e\in\hlfcfib}e(h)=\sum_{e\in\mathcal F}\hlfcloc_a(A_e)$.
\end{enumerate}
\formallabel{UpperTheory.Gal.productCup_of_system.yMain_of}
\end{lemma}

\begin{proof}
By \cref{pf:hlf-chart:crt} choose $h\in B'_{\{0\}}$ with $e(h)=\hlfcloc_a(A_e)$ for $e\in\mathcal F$ and $e(h)=0$ for $e\in\hlfcfib\smallsetminus\mathcal F$;
this gives (c). For (b): $\omega'_0\in\Omega^d(B'_{[d]})$ since the $t_l$ are units in $B'_{[d]}$, and a cochain concentrated at $[d]$ in degree $d$ is
a cocycle (there are no sets of cardinality $d+2$).

For (a), let $c_0=c_{\{0\}\sqcup\{\bar 0\}\sqcup D_0}$ (the $D_0$-minors, with no index from $\kappa$). Its column lies in $u\,M^Z_{\{0\},\{0\}}$ with
$u=(X_0Y_0)^d$ (\cref{pf:hlf-chart:sysbasic}(a)), so $c_0=u\,m'$ with $m'\in\Omega^{d-a}(B_{\{0\}})\otimes\Omega^a(B_{\{0\}})$. As the restriction of $c$
to the minors is diagonal, $c_0$ acts as $(-1)^d\kappa_{D_0}$; hence $m=(-1)^dm'=(-1)^du^{-1}c_0$ acts as
$\hlfcmult((-1)^du^{-1})(-1)^d\kappa_{D_0}=x_0^{-2d}\kappa_{D_0}=\prod_{l\ge1}t_l$ (the computation of $\kappa_{D_0}$ above; kernels which act form
a module, \lref{UpperTheory.Gal.productCup_of_system.actsLin_of}). By \cref{pf:hlf-chart:valg}, $\gamma(m)=s\,dt_\star$ with
$s(e)=[X^a]\det(1+XA_e)=e(h)\det(1-A_e)$ for $e\in\mathcal F$, and therefore
\[
  \gamma(c_0)=g\bigl((-1)^du\bigr)\gamma(m)=(-1)^d\bigl(x_0j_2(x_0)\bigr)^d s\,dt_\star .
\]
By \cref{pf:hlf-chart:gmap}, $c'=\gamma(c)$ is a cocycle of degree $d+2$ of $\mathrm{Tot}_{T_Y}$ for the fixed-point data, with component $\gamma(c_0)$ at
$\{0\}\sqcup\{\bar0\}\sqcup D_0$. Let $c_1$ be the cocycle of \cref{pf:hlf-chart:ybig}(a). By \cref{pf:hlf-chart:ybig}(c),
$\mathrm{top}(c')-\mathrm{top}(c_1)$ is a coboundary of $\mathrm{Mix}(T_Y)$. The restriction of $c_1$ to the origin data is a cocycle
(\cref{pf:hlf-chart:sysbasic}(d)) with the same top part, the component $(-1)^dx_0^dh\,dt_\star$ at $\{0\}\sqcup\{\bar0\}\sqcup\{1..d\}$ and vanishing
$(\{i\},\{\bar j\})$-columns for $(i,j)\ne(0,0)$; by \cref{pf:hlf-chart:ybig}(b), $\mathrm{top}(c_1)-\mathrm{edge}_1(h\omega'_0)$ is a coboundary. Adding,
and using $\mathrm{top}(\gamma(c))=\gamma(\mathrm{top}(c))$, gives (a).
\end{proof}

\subsection{Proof of the core statement}\label{pf:hlf-chart:ss:core}

\begin{proof}[Proof of \cref{pf:hlf-chart:core}]
Fix a top coefficient $\tau$ for $K$ (\cref{pf:hlf-chart:dlog}) and a residue functional $\rho$ (\cref{pf:hlf-chart:rho}); let $\theta$ be the trace
functional of $K$ and $\langle\ ,\ \rangle$ the pairing \eqref{eq:pf:hlf-chart:pair}.

\emph{Step 1: the diagonal.} Choose diagonal Koszul data $(n,z)$ (\cref{pf:hlf-chart:zfam}) and a cocycle $c\in Z^{d+2}(\mathrm{Tot}_{T_Z}(\mu\sqcup z))$
whose restriction $c^\mu$ to the minors is diagonal (\cref{pf:hlf-chart:zex}); $c^\mu$ is a cocycle (\cref{pf:hlf-chart:sysbasic}(d)) and
$\mathrm{top}(c^\mu)=\mathrm{top}(c)$. By \cref{pf:hlf-chart:zid} there are triples $(p_k,e_k,f_k)_{1\le k\le M}$ with $p_k\le d$,
$e_k\in Z^{d-p_k}(\hlfcC^\bullet(B_\bullet,d-a))$, $f_k\in Z^{p_k}(\hlfcC^\bullet(B_\bullet,a))$, satisfying \eqref{eq:pf:hlf-chart:topdec} and
\eqref{eq:pf:hlf-chart:dualdec}.

\emph{Step 2: the trace as a sum of pairings.} By \ref{pf:hlf-chart:H10} applied to $f_k$ there are $\tilde b_k\in Z^{p_k}(\hlfcC^\bullet(B'_{\bullet\bullet},a))$
with $(\tilde b_k)_{\emptyset\sqcup\bar J}=j_2^*((f_k)_J)$ for all $J$, such that $g_k=(\Tr(\tilde b_k)_{I\sqcup\emptyset})_I$ is a $p_k$-cocycle of
$\hlfcC^\bullet(B_\bullet,a)$ with $u_{p_k}[f_k]=[g_k]$. By \cref{pf:hlf-chart:alttrace},
\begin{equation}\label{eq:pf:hlf-chart:step2}
  \sum_{q=0}^d(-1)^q\tr(u_q)=\sum_k(-1)^{p_k}\langle e_k,g_k\rangle .
\end{equation}

\emph{Step 3: a cocycle on $Y$.} Put $a_k=(j_1^*(e_k)_I)_I$ and $b_k=((\tilde b_k)_{I\sqcup\emptyset})_I$. Since $j_1(B_I)\subseteq B'_I$
(\cref{pf:hlf-chart:rybasic}) and $j_1^*$ commutes with $\delta$, $a_k\in Z^{d-p_k}(\hlfcC^\bullet(B'_\bullet,d-a))$; $b_k$ is the restriction of the
cocycle $\tilde b_k$ to the index sets contained in the first copy of $[d]$, hence $b_k\in Z^{p_k}(\hlfcC^\bullet(B'_\bullet,a))$ ($B'_{I,\emptyset}=B'_I$). Let
\[
  \Lambda=\sum_k(-1)^{p_k}\,a_k\cup_\wedge b_k ,
\]
a $d$-cocycle of $\hlfcC^\bullet(B'_\bullet,d)$ (\cref{pf:hlf-chart:lcup}, or the Leibniz rule; $\Omega^{d-a}(B'_{I'})\wedge\Omega^a(B'_{I''})\subseteq\Omega^d(B'_I)$ for
$I',I''\subseteq I$ by \cref{pf:hlf-chart:rybasic,pf:hlf-chart:omloc}).

\emph{Step 4: comparison with the pulled-back diagonal.} Apply \cref{pf:hlf-chart:lcup} with $\mu=\wedge$, $A_I=\Omega^{d-a}(B'_I)$,
$G_{I\sqcup\bar J}=\Omega^a(B'_{I,J})$ and $H_{I,J}=\Omega^d(B'_{I,J})$; its hypothesis holds by the monotonicity of $B'_{I,J}$
(\cref{pf:hlf-chart:rybasic}(a)) and \cref{pf:hlf-chart:omloc}. With $a=a_k$, $\gamma=\tilde b_k$ we get
\[
  \mathrm{join}_\wedge(a_k,\tilde b_k^{(2)})-(-1)^{p_k}\,\mathrm{edge}_1(a_k\cup_\wedge b_k)\in B^{d+2}(\mathrm{Mix}(T_Y)).
\]
Here $\mathrm{join}_\wedge(a_k,\tilde b_k^{(2)})_{I\sqcup\bar J}=j_1^*(e_k)_I\wedge j_2^*(f_k)_J=\gamma\bigl((e_k\boxtimes f_k)_{I\sqcup\bar J}\bigr)$. By
\eqref{eq:pf:hlf-chart:topdec} and \cref{pf:hlf-chart:gmap}, $\gamma(\mathrm{top}(c))-\sum_k\gamma(e_k\boxtimes f_k)$ is a coboundary of
$\mathrm{Mix}(T_Y)$. Summing,
\[
  \gamma(\mathrm{top}(c))-\mathrm{edge}_1(\Lambda)\in B^{d+2}(\mathrm{Mix}(T_Y)).
\]

\emph{Step 5: localization.} Let $h\in B'_{\{0\}}$ be as in \cref{pf:hlf-chart:ymain}. Subtracting \cref{pf:hlf-chart:ymain}(a) from the last relation,
$\mathrm{edge}_1(\Lambda-h\omega'_0)\in B^{d+2}(\mathrm{Mix}(T_Y))$, where $\Lambda-h\omega'_0$ is a $d$-cocycle by \cref{pf:hlf-chart:ymain}(b). By
\cref{pf:hlf-chart:edgeb}, $\Lambda-h\omega'_0\in B^d(\hlfcC^\bullet(B'_\bullet,d))$. By \cref{pf:hlf-chart:trtop}, \cref{pf:hlf-chart:rt1} (for $L=K'$,
using \ref{pf:hlf-chart:H1} and \ref{pf:hlf-chart:H3}) and \cref{pf:hlf-chart:ymain}(c),
\begin{equation}\label{eq:pf:hlf-chart:step5}
  \theta\bigl(\Tr(\Lambda_{[d]})\bigr)=\theta\bigl(\Tr(h\omega'_0)\bigr)=\rho\bigl(\Tr_{K'/F}h\bigr)=\sum_{e\in\hlfcfib}e(h)=\sum_{e\in\mathcal F}\hlfcloc_a(A_e).
\end{equation}

\emph{Step 6: the pairings.} For each $k$, by the projection formula (\cref{pf:hlf-chart:trdr}),
\[
  \Tr\bigl((a_k\cup_\wedge b_k)_{[d]}\bigr)=\sum_j\Tr\bigl(j_1^*(e_k)_{\{0..j\}}\wedge(\tilde b_k)_{\{j..d\}\sqcup\emptyset}\bigr)
  =\sum_j(e_k)_{\{0..j\}}\wedge(g_k)_{\{j..d\}},
\]
so $\theta(\Tr((a_k\cup b_k)_{[d]}))=\langle e_k,g_k\rangle$ and $\theta(\Tr(\Lambda_{[d]}))=\sum_k(-1)^{p_k}\langle e_k,g_k\rangle$. Comparing with
\eqref{eq:pf:hlf-chart:step2} and \eqref{eq:pf:hlf-chart:step5} gives $\sum_{q=0}^d(-1)^q\tr(u_q)=\sum_{e\in\mathcal F}\hlfcloc_a(A_e)$.
\end{proof}

\begin{remark}[Scope]\label{pf:hlf-chart:rem:scope}
The proof uses the operators $u_q$ only through the lifts of the finitely many cocycles $f_k$, and it shows that the spaces $\hlfcH^q(B_\bullet,a)$
are finite-dimensional (\cref{pf:hlf-chart:alttrace}), which is why \ref{pf:hlf-chart:H11} is not needed. The hypothesis \ref{pf:hlf-chart:H7} is
not used.
\end{remark}
